\documentclass[11pt,oneside]{article}

\usepackage[a4paper,margin=25mm,headheight=15pt]{geometry}
\usepackage[T1]{fontenc}
\usepackage[utf8]{inputenc}
\IfFileExists{libertinus.sty}{\usepackage{libertinus}}{\usepackage{lmodern}}
\usepackage{microtype}
\usepackage{amsmath,amssymb,amsthm,mathtools,mathrsfs,bm}
\usepackage{booktabs,longtable,tabularx,array,multirow}
\usepackage{graphicx,float,caption,subcaption}
\usepackage{xcolor}
\usepackage{enumitem}
\usepackage{fancyhdr}
\usepackage{titlesec}
\usepackage{xurl}
\usepackage{listings}
\usepackage{hyperref}
\usepackage[nameinlink,noabbrev]{cleveref}

\definecolor{linkblue}{RGB}{24,72,124}
\definecolor{deepgreen}{RGB}{23,102,69}
\definecolor{deepred}{RGB}{140,42,42}
\definecolor{shadegray}{RGB}{246,247,249}
\definecolor{rulegray}{RGB}{195,201,208}
\definecolor{deepviolet}{RGB}{86,53,120}
\hypersetup{
  colorlinks=true,
  linkcolor=linkblue,
  citecolor=deepgreen,
  urlcolor=linkblue,
  pdftitle={Signed and Reciprocal q-Fabius Frontiers},
  pdfauthor={OpenAI GPT-5.6 Pro},
  pdfsubject={Negative parameters, reciprocal Laplace duality, two-nome Pochhammer-Prouhet partition functions, and inverse-geometric endpoint lattices},
  pdfkeywords={Fabius function, Rvachev up-function, Thue-Morse, q-binomial, q-Pochhammer, negative q, reciprocal q, Bell polynomials, Bernoulli polynomials, Lambert W}
}

\pagestyle{fancy}
\fancyhf{}
\fancyhead[L]{\small Signed and reciprocal $q$-Fabius frontiers}
\fancyhead[R]{\small\nouppercase{\leftmark}}
\fancyfoot[C]{\thepage}
\renewcommand{\headrulewidth}{0.4pt}

\titleformat{\section}{\Large\bfseries\color{linkblue}}{\thesection}{0.72em}{}
\titleformat{\subsection}{\large\bfseries\color{linkblue}}{\thesubsection}{0.65em}{}
\titleformat{\subsubsection}{\normalsize\bfseries\color{deepgreen}}{\thesubsubsection}{0.58em}{}
\setcounter{tocdepth}{2}
\setcounter{secnumdepth}{3}
\numberwithin{equation}{section}
\allowdisplaybreaks[3]
\setlength{\emergencystretch}{3em}
\setlist[itemize]{leftmargin=2em,itemsep=0.25em,topsep=0.45em}
\setlist[enumerate]{leftmargin=2.3em,itemsep=0.25em,topsep=0.45em}

\newtheorem{theorem}{Theorem}[section]
\newtheorem{proposition}[theorem]{Proposition}
\newtheorem{lemma}[theorem]{Lemma}
\newtheorem{corollary}[theorem]{Corollary}
\newtheorem{conjecture}[theorem]{Conjecture}
\newtheorem{problem}[theorem]{Open problem}
\theoremstyle{definition}
\newtheorem{definition}[theorem]{Definition}
\newtheorem{algorithm}[theorem]{Algorithm}
\newtheorem{example}[theorem]{Example}
\theoremstyle{remark}
\newtheorem{remark}[theorem]{Remark}
\newtheorem{warning}[theorem]{Warning}

\newcommand{\R}{\mathbb R}
\newcommand{\C}{\mathbb C}
\newcommand{\Q}{\mathbb Q}
\newcommand{\N}{\mathbb N}
\newcommand{\Z}{\mathbb Z}
\newcommand{\Fq}{F_q}
\newcommand{\Gq}{G_q}
\newcommand{\Up}{\operatorname{up}}
\newcommand{\sinc}{\operatorname{sinc}}
\newcommand{\sgn}{\operatorname{sgn}}
\newcommand{\supp}{\operatorname{supp}}
\newcommand{\ord}{\operatorname{ord}}
\newcommand{\Var}{\operatorname{Var}}
\newcommand{\Cov}{\operatorname{Cov}}
\newcommand{\Law}{\operatorname{Law}}
\newcommand{\Prob}{\mathbb P}
\newcommand{\E}{\mathbb E}
\newcommand{\dd}{\,\mathrm d}
\newcommand{\e}{\mathrm e}
\newcommand{\ii}{\mathrm i}
\newcommand{\bigO}{\mathcal O}
\newcommand{\smallo}{o}
\newcommand{\qpoch}[3]{\left(#1;#2\right)_{#3}}
\newcommand{\qbinom}[3]{\genfrac{[}{]}{0pt}{}{#1}{#2}_{#3}}
\newcommand{\Bell}{\mathcal B}
\newcommand{\wt}{\operatorname{wt}_2}
\newcommand{\vtwo}{\nu_2}
\newcommand{\repo}{\textsc{ProveIt}}
\newcommand{\proved}{\textcolor{deepgreen}{\textbf{proved here}}}
\newcommand{\derived}{\textcolor{linkblue}{\textbf{derived from corpus infrastructure}}}
\newcommand{\openstatus}{\textcolor{deepred}{\textbf{conjectural/open}}}
\newcommand{\auditnew}{\textcolor{deepviolet}{\textbf{new relative to the pinned snapshot}}}
\newcommand{\abs}[1]{\left\lvert#1\right\rvert}
\newcommand{\norm}[1]{\left\lVert#1\right\rVert}
\newcommand{\floor}[1]{\left\lfloor#1\right\rfloor}
\newcommand{\ceil}[1]{\left\lceil#1\right\rceil}
\newcommand{\pospart}[1]{\left(#1\right)_{+}}
\newcommand{\Br}{\mathfrak B}
\newcommand{\Zpart}{\mathscr Z}
\newcommand{\Dcomb}{\Delta}
\newcommand{\Gr}{\operatorname{Gr}}

\newenvironment{statusbox}[1]{%
  \par\medskip\noindent\begin{center}\begin{minipage}{0.94\textwidth}%
  \color{linkblue}\hrule\vspace{0.65em}%
  \color{black}\textbf{\color{linkblue}#1.}\ }{%
  \par\vspace{0.65em}\color{linkblue}\hrule
  \end{minipage}\end{center}\medskip}

\newenvironment{resultbox}[1]{%
  \par\medskip\noindent\begin{center}\begin{minipage}{0.94\textwidth}%
  \color{deepgreen}\hrule\vspace{0.65em}%
  \color{black}\textbf{\color{deepgreen}#1.}\ }{%
  \par\vspace{0.65em}\color{deepgreen}\hrule
  \end{minipage}\end{center}\medskip}

\lstdefinestyle{python}{
  language=Python,
  basicstyle=\ttfamily\small,
  backgroundcolor=\color{shadegray},
  frame=single,
  rulecolor=\color{rulegray},
  columns=fullflexible,
  breaklines=true,
  keepspaces=true,
  showstringspaces=false,
  keywordstyle=\color{linkblue}\bfseries,
  commentstyle=\color{deepgreen},
  stringstyle=\color{deepred}
}

\title{\bfseries Signed and Reciprocal $q$-Fabius Frontiers\\[0.3em]
\large Negative parameters, reciprocal Laplace duality, two-nome\\
Pochhammer--Prouhet partition functions, and inverse-geometric endpoint lattices}
\author{A frontier research report based on the recursive \texttt{*.tex} corpus under\\
\small\path{Analysis/FabiusFunction/docs}}
\date{Pinned repository snapshot: 28 August 2026\\
\small commit \texttt{89b1a5f915cad065a42ee48f758eb2cd2d9beb94}}

\begin{document}
\maketitle

\begin{abstract}
The existing \repo{} corpus develops the classical Fabius distribution and Rvachev
up-function at dyadic scale, exact dyadic arithmetic, the Thue--Morse finite-block
calculus, half-base Gaussian coefficients, $q$-Pochhammer products, finite sinc
products, Bell--Bernoulli endpoint expansions, Lambert-$W_{-1}$ saddle coordinates,
and several $q=1/4$ Richardson schemes.  The purpose of this report is not to
repackage that material.  It studies instead the signed and reciprocal parameter
orbit
\[
 \left\{q,-q,q^{-1},-q^{-1}\right\}
\]
after separating three roles that are frequently denoted by the same letter:
the geometric ratio of the random series, the Gaussian-binomial nome, and the
digit-position weight in a Thue--Morse projection.

For $|q|<1$ define the normalized geometric-uniform variable
\[
 Y_q=(1-q)\sum_{j\ge0}q^jU_j,
 \qquad U_j\sim\operatorname{Unif}[0,1]
\]
independently, and let $M_q(t)=\E\e^{tY_q}$.  The case $q=1/2$ is exactly the
Fabius distribution.  We extend $M_q$ to $|q|>1$ as a normalized meromorphic
moment germ satisfying the same Mahler equation.  The resulting system obeys
\[
 M_q(t)M_{1/q}(-t)=1,
 \qquad
 Y_{-a}\stackrel d=\frac{1+a}{1-a}Y_a-\frac{a}{1-a}
 \quad(0<a<1).
\]
Thus the requested values split into two Klein-four orbits,
$\{\pm\frac12,\pm2\}$ and $\{\pm\frac14,\pm4\}$, and every member is obtained
from one positive contraction by affine transport and reciprocal inversion.
The spectral product has nome $q^2$: sign change rescales the centered germ,
whereas inversion exchanges its zeros and poles.

Several further consequences are proved.  First, the reciprocal centered germ
is a genuine characteristic function: for $|q|>1$ it is the Fourier transform
of an explicit infinite convolution of Laplace laws.  Hence the formal moments
are moments of a positive measure on the vertical line $\Re z=1/2$, their
Hankel determinants have exact sign $(-1)^{\binom n2}$, and their monic
orthogonal polynomials have all zeros on that line.  Second, the inverse
averaging operator produces a $q$-Fabius--Bernoulli Appell family interpolating
ordinary Bernoulli polynomials at $q=0$ and centered monomials as $q\to1^-$.  A
new polynomial
\[
 \mathcal P_n(q)=\frac{(q;q)_n}{(1-q)^n}\frac{\E Y_q^n}{n!}
\]
packages the moments and yields an exact $q$-Pochhammer formula for the
geometric endpoint lattice.

The Thue--Morse and Gaussian layers are unified by the two-nome binary-cube
partition function
\[
 \Zpart_N(u,t;\rho,\sigma)
 =\prod_{j=0}^{N-1}\left(1-u\rho^j\e^{t\sigma^j}\right).
\]
Its $t=0$ face is finite $q$-Pochhammer/Gaussian-binomial calculus; its
$(u,\rho)=(1,1)$ face is a geometric Prouhet generator.  From it we obtain an
exact $q$-Prouhet moment-transfer theorem, signed position-energy layers,
simultaneous multisection in both nomes, and a finite-field/unitary square for
$q=\pm1/2,\pm2,\pm1/4,\pm4$.  The associated box-spline derivative comb is the
full dyadic Thue--Morse grid at $q=1/2$ but a signed Cantor skeleton of Hausdorff
dimension $1/2$ at $q=1/4$.

Finally, for $0<q\le1/2$ we prove the inverse-geometric endpoint identity
\[
 G_q(q^n)=
 \frac{q^{\binom{n+1}{2}}}{(1-q)^n}\frac{\E Y_q^n}{n!}
 =\frac{q^{\binom{n+1}{2}}}{(q;q)_n}\mathcal P_n(q),
\]
including all derivative jets there.  It gives an exact inverse-quartic lattice
for $q=1/4$ and yields the uniform two-term endpoint and inverse asymptotics
\[
 \log G_q(x)=-\frac{\log^2(1/x)}{2\log(1/q)}
 -\frac{\log(1/x)\log\log(1/x)}{\log(1/q)}+O(\log(1/x)),
\]
\[
 \log\frac1{G_q^{-1}(y)}
 =\sqrt{2\log(1/q)\log(1/y)}-\frac12\log\log(1/y)+O(1).
\]
A generalized Lambert-$W_{-1}$ all-orders program, natural-boundary
conjectures, arithmetic questions for $\mathcal P_n$, and a Lean formalization
roadmap conclude the report.  A fully commented Python program verifies the
principal orbit, Prouhet, endpoint, determinant, and moment-polynomial identities
at the eight requested parameter values and supplies the figures and CSV tables
in the accompanying archive.
\end{abstract}

\tableofcontents
\clearpage

\section{Audit boundary, terminology, and principal results}
\label{sec:audit}

\subsection{Pinned corpus and novelty discipline}

The source boundary for this report is commit
\texttt{89b1a5f915cad065a42ee48f758eb2cd2d9beb94}.  The repository was changing
on the audit date, so every statement of novelty is relative to this exact
snapshot.  The audit followed the repository's own stratification:
\begin{itemize}
\item the proof-backed primary exposition and its Lean coverage map;
\item the canonical non-formalized frontier volume;
\item the thematic synthesis volumes for Thue--Morse, exponent sequences and
      $q$-series, spectra and arithmetic, representations, inverse sampling,
      and integration/transforms;
\item archived $q$-calculus and dyadic-connection reports, treated as
      mathematical lineage rather than independent current claims;
\item targeted code and formula searches across the recursive documentation
      tree for the identities proposed below.
\end{itemize}

The most important prior-art boundary is the consolidated volume
\path{non-formalized-research-frontiers/drafts/exponents-and-q-series/}
\path{Exponents_and_q_Series_Frontiers/Exponents_and_q_Series_Frontiers.tex}.
It already contains arbitrary nonnegative exponent-sequence sinc products,
zero multiplicities, spectral zeta and cumulant dictionaries, Gaussian
$q$-binomial tomography, the Pascal--Rvachev hierarchy, $q=1/4$ Richardson
acceleration, finite-product spline approximants, sharp transport laws, and
Mellin--Lambert endpoint mechanisms.  The Thue--Morse atlas already contains
the standard product, Prouhet cancellation, Walsh/Fourier formulae, automatic
and rarefied sums, and the bridge to the Fabius/Rvachev system.  None of those
is claimed as new here.

Targeted searches in the pinned tree found no exact counterpart of the
following package:
\begin{enumerate}
\item the normalized signed/reciprocal moment germ on the full orbit
      $\{q,-q,q^{-1},-q^{-1}\}$ and its Cayley-coordinate action;
\item the affine conjugacy of negative geometric-uniform laws;
\item the reciprocal Laplace-convolution representation, exact Hankel
      signature, and vertical-line orthogonality;
\item the two-nome partition function combining Gaussian position energy and
      geometric Prouhet projection;
\item the all-$q$ finite Prouhet moment-transfer identity;
\item the inverse-geometric endpoint formula, its derivative jets, and its
      $q=1/4$ specialization;
\item the $q$-Fabius--Bernoulli inverse-averaging Appell family;
\item the moment polynomial $P_n(q)$ and its endpoint $q$-Pochhammer form.
\end{enumerate}
External searches located the negative-$q$ Gaussian-binomial interpretation of
Fu, Reiner, Stanton, and Thiem, but not the Fabius/Rvachev reciprocal and
endpoint structures assembled here.  Accordingly, \auditnew{} means only
``not found in the pinned corpus after targeted audit and not found as a named
package in the targeted literature search.''  It is not a claim of exhaustive
worldwide priority.

\begin{statusbox}{Proof-status convention}
Every theorem marked \proved{} has a proof in this manuscript.  The label
\derived{} means that the argument uses an established object or theorem from
the audited corpus and adds a new deduction.  Statements marked
\openstatus{} are explicitly conjectural.  Numerical evidence never replaces
a proof; its role is reproducibility, error detection, and conjecture triage.
\end{statusbox}

\subsection{Three nomes that should not be conflated}

The letter $q$ occurs in several mathematically different roles.  We keep them
separate whenever two of them coexist:
\begin{center}
\begin{tabularx}{0.95\textwidth}{>{\bfseries}l l X}
\toprule
symbol & role & typical object \\
\midrule
$q$ & geometric ratio & $Y_q=(1-q)\sum q^jU_j$, $M_q$, $\Phi_q$ \\
$\rho$ & position-energy/Gaussian nome & $(u;\rho)_N$, $\qbinom{N}{k}{\rho}$ \\
$\sigma$ & geometric digit projection & $A_{\sigma,N}(\epsilon)=\sum\epsilon_j\sigma^j$ \\
$Q$ & finite-field cardinality & $\mathbb F_Q$, Hermitian space over $\mathbb F_{Q^2}$ \\
$q^2$ & spectral Pochhammer nome & zeros/poles of the centered sinc/hyperbolic product \\
\bottomrule
\end{tabularx}
\end{center}
The classical Fabius case sets $q=1/2$.  A finite Gaussian stencil derived from
that case often also has nome $\rho=1/2$, but this equality is contingent, not
structural.  The master partition function in \cref{sec:two-nome} allows
$\rho$ and $\sigma$ to vary independently and thereby exposes which statement
belongs to which layer.

\subsection{The two requested Klein-four orbits}

The transformations $s(q)=-q$ and $r(q)=q^{-1}$ commute and generate a
Klein four-group away from $q=0$.  The values requested in the problem are
exactly two orbits:
\[
 \mathcal O_2=\left\{\frac12,-\frac12,2,-2\right\},
 \qquad
 \mathcal O_4=\left\{\frac14,-\frac14,4,-4\right\}.
\]
The following table previews their formal means and variances.  The word
``formal'' is needed only outside the unit disk; inside it these are ordinary
probability moments.

\begin{table}[H]
\centering
\caption{The requested signed/reciprocal parameter orbits.}
\label{tab:orbit-variance}
\begin{tabular}{c c c c}
\toprule
$q$ & geometric status & mean & variance/formal second cumulant \\
\midrule
$1/2$ & Fabius probability law & $1/2$ & $1/36$ \\
$-1/2$ & affine Fabius law on $[-1,2]$ & $1/2$ & $1/4$ \\
$2$ & reciprocal moment germ & $1/2$ & $-1/36$ \\
$-2$ & reciprocal signed germ & $1/2$ & $-1/4$ \\
\addlinespace
$1/4$ & quartic geometric probability law & $1/2$ & $1/20$ \\
$-1/4$ & affine quartic law on $[-1/3,4/3]$ & $1/2$ & $5/36$ \\
$4$ & reciprocal quartic germ & $1/2$ & $-1/20$ \\
$-4$ & reciprocal signed quartic germ & $1/2$ & $-5/36$ \\
\bottomrule
\end{tabular}
\end{table}

\subsection{Executive theorem map}

The report has four mutually reinforcing layers.
\begin{enumerate}
\item \textbf{Moment/Mahler layer.}  A single functional equation defines a
      probability MGF for $|q|<1$ and a reciprocal germ for $|q|>1$.  Bell and
      Bernoulli polynomials give every coefficient.
\item \textbf{Spectral layer.}  The centered transform is an infinite
      $q^2$-Pochhammer product.  Sign change is an affine rescaling; inversion
      exchanges zeros and poles and produces a positive Laplace spectral dual.
\item \textbf{Binary-cube layer.}  A two-nome partition function contains
      Gaussian coefficients, signed position-energy layers, Thue--Morse
      Prouhet cancellation, and multisection as different specializations.
\item \textbf{Endpoint layer.}  A simplex-volume argument transfers the
      moment system to an exact geometric endpoint lattice and its derivative
      jets, with a uniform first asymptotic inversion and a conjectural
      all-orders Lambert-$W_{-1}$ continuation.
\end{enumerate}

\section{The normalized geometric-uniform law and its reciprocal germ}
\label{sec:master-germ}

\subsection{The positive and signed probability regimes}

Let $U_0,U_1,\ldots$ be independent $\operatorname{Unif}[0,1]$ random
variables.  For a real or complex parameter $q$ with $|q|<1$, define
\begin{equation}
 Y_q=(1-q)\sum_{j=0}^{\infty}q^jU_j.
 \label{eq:Yq}
\end{equation}
The series converges absolutely, and its weights sum to one.  Consequently
\begin{equation}
 \E Y_q=\frac12
 \label{eq:Yq-mean}
\end{equation}
for every real $|q|<1$, including negative $q$.

For $0\le q<1$, all weights are nonnegative and the support is $[0,1]$.  The
case $q=1/2$ is the random-series representation of the classical bounded
Fabius function.  We write
\begin{equation}
 G_q(x)=\Prob(Y_q\le x)
 \label{eq:Gq}
\end{equation}
for the CDF in the positive regime.  The existing Lean API in
\path{GeometricUniformLaw.lean} formalizes the normalized weights, the affine
self-similarity, reflection, absolute continuity, null singletons, and support.

For $q=-a$, $0<a<1$, positive and negative coefficient intervals add to
\[
 \sum_{j\ \mathrm{even}}(1+a)a^j=\frac1{1-a},
 \qquad
 -\sum_{j\ \mathrm{odd}}(1+a)a^j=-\frac a{1-a}.
\]
Hence
\begin{equation}
 \supp(Y_{-a})=\left[-\frac a{1-a},\frac1{1-a}\right].
 \label{eq:negative-support}
\end{equation}
The support remains an interval, not a Cantor set, because each coordinate is
itself an interval variable.

\subsection{The universal Mahler equation}

Put
\begin{equation}
 E(z)=\frac{\e^z-1}{z},\qquad E(0)=1,
 \qquad
 S(z)=\e^{-z/2}E(z)=\frac{\sinh(z/2)}{z/2}.
 \label{eq:E-S}
\end{equation}
For $|q|<1$, independence gives the entire moment generating function
\begin{equation}
 M_q(t)=\E\e^{tY_q}
 =\prod_{j=0}^{\infty}E\!\left((1-q)q^jt\right).
 \label{eq:Mq-inside}
\end{equation}
Splitting the factor $j=0$ gives
\begin{equation}
 \boxed{M_q(t)=E((1-q)t)M_q(qt).}
 \label{eq:Mahler}
\end{equation}
This is the normalized geometric analogue of the classical Fabius moment
Mahler equation.

\begin{definition}[Reciprocal normalized germ]
For $|q|>1$, define near $t=0$
\begin{equation}
 M_q(t)=
 \left[
  \prod_{j=1}^{\infty}E\!\left((1-q)q^{-j}t\right)
 \right]^{-1}.
 \label{eq:Mq-outside}
\end{equation}
It is meromorphic globally and analytic at the origin.  Its Taylor
coefficients are called the \emph{reciprocal $q$-Fabius moments}.
\end{definition}

The indexing in \eqref{eq:Mq-outside} is forced by normalization.  Reindexing
immediately proves that the same equation \eqref{eq:Mahler} holds for
$|q|>1$.  Thus the unit circle, rather than the sign of $q$, is the natural
dividing curve between the probability and reciprocal-germ regimes.

\begin{theorem}[Uniqueness of the normalized formal solution]
\label{thm:formal-unique}
Let $q\in\C$ satisfy $q^n\ne1$ for every $n\ge1$.  There is a unique formal
power series $M_q(t)\in\C[[t]]$ with $M_q(0)=1$ satisfying
\eqref{eq:Mahler}.  It is represented by \eqref{eq:Mq-inside} when $|q|<1$
and by \eqref{eq:Mq-outside} when $|q|>1$.  This statement is \proved.
\end{theorem}

\begin{proof}
Write $M_q(t)=\sum_{n\ge0}m_n(q)t^n/n!$.  Comparing the coefficient of $t^n$
in \eqref{eq:Mahler} gives
\begin{equation}
 (1-q^n)m_n(q)
 =\sum_{k=0}^{n-1}\binom nk
   \frac{(1-q)^{n-k}q^k}{n-k+1}m_k(q).
 \label{eq:moment-recurrence}
\end{equation}
The coefficient $1-q^n$ is nonzero, so $m_n$ is determined recursively by
$m_0=1$.  The convergent products satisfy the same equation and normalization,
therefore realize the unique formal solution.
\end{proof}

For exponential coefficients $a_n(q)=m_n(q)/n!$, the recurrence is
\begin{equation}
 (1-q^n)a_n(q)=
 \sum_{r=1}^{n}
 \frac{(1-q)^rq^{n-r}}{(r+1)!}a_{n-r}(q).
 \label{eq:an-recurrence}
\end{equation}
It specializes at $q=1/2$ to the half-moment recurrence used throughout the
classical Fabius literature.

\subsection{Reflection and centered normalization}

Since $E(-z)=\e^{-z}E(z)$ and the normalized weights sum to one, both convergent
products and the formal solution satisfy
\begin{equation}
 M_q(-t)=\e^{-t}M_q(t).
 \label{eq:M-reflection}
\end{equation}
Define the centered germ
\begin{equation}
 C_q(t)=\e^{-t/2}M_q(t).
 \label{eq:Cq}
\end{equation}
Then
\begin{equation}
 C_q(-t)=C_q(t),
 \label{eq:C-even}
\end{equation}
so every centered odd cumulant vanishes.  In the probability regime this is
ordinary reflection of $Y_q$ about $1/2$; outside the unit disk it remains an
exact algebraic symmetry of the germ.

\section{Bernoulli cumulants and Bell-polynomial moments}
\label{sec:cumulants}

\subsection[All cumulants in one rational function of q]{All cumulants in one rational function of $q$}

The expansion
\begin{equation}
 \log E(z)=\frac z2+
 \sum_{n=2}^{\infty}\frac{B_n^+}{n}\frac{z^n}{n!}
 \label{eq:logE}
\end{equation}
uses the Bernoulli convention $B_1^+=+1/2$; for $n\ne1$ these are the ordinary
Bernoulli numbers.  Summing over the geometric scales gives the following
closed form.

\begin{theorem}[Signed/reciprocal cumulant dictionary]
\label{thm:cumulants}
For every nonsingular $q$,
\begin{equation}
 \boxed{
 \kappa_1(q)=\frac12,
 \qquad
 \kappa_n(q)=\frac{B_n^+}{n}\frac{(1-q)^n}{1-q^n}
 \quad(n\ge2).}
 \label{eq:cumulants}
\end{equation}
In particular, $\kappa_{2m+1}=0$ for $m\ge1$ and
\begin{equation}
 \kappa_2(q)=\frac1{12}\frac{1-q}{1+q}.
 \label{eq:variance}
\end{equation}
The theorem is \proved.
\end{theorem}

\begin{proof}
For $|q|<1$, insert \eqref{eq:logE} into \eqref{eq:Mq-inside} and sum
$\sum_{j\ge0}q^{jn}=(1-q^n)^{-1}$.  Both sides of \eqref{eq:cumulants} are
rational functions of $q$ coefficient by coefficient and satisfy the unique
formal recurrence of \cref{thm:formal-unique}; equivalently one may repeat the
calculation in the reciprocal product.  The variance formula follows from
$B_2=1/6$.
\end{proof}

The simple Möbius function $(1-q)/(1+q)$ already predicts the sign pattern in
\cref{tab:orbit-variance}: inversion changes the sign of every even centered
cumulant,
\begin{equation}
 \kappa_{2m}(1/q)=-\kappa_{2m}(q).
 \label{eq:cumulant-inversion}
\end{equation}
Sign change rescales them by an even power, as proved in
\cref{sec:klein-four}.

\subsection{Complete Bell-polynomial reconstruction}

Let $\Bell_n(x_1,\ldots,x_n)$ denote the complete exponential Bell polynomial,
characterized by
\[
 \exp\!\left(\sum_{j\ge1}x_j\frac{t^j}{j!}\right)
 =\sum_{n\ge0}\Bell_n(x_1,\ldots,x_n)\frac{t^n}{n!}.
\]
Then
\begin{equation}
 \boxed{
 m_n(q)=\Bell_n\left(
 \frac12,\kappa_2(q),0,\kappa_4(q),0,\ldots
 \right).}
 \label{eq:Bell-moments}
\end{equation}
This makes all requested parameter values exactly computable over $\Q$.  The
first moments are
\begin{align}
 m_0&=1,
 &m_1&=\frac12,\nonumber\\
 m_2&=\frac14+\kappa_2,
 &m_3&=\frac18+\frac32\kappa_2,\nonumber\\
 m_4&=\frac1{16}+\frac32\kappa_2+3\kappa_2^2+\kappa_4.
 \label{eq:first-moments-cumulants}
\end{align}
The accompanying CSV file \path{data/q_orbit_moments.csv} lists $m_0$ through
$m_6$ for all eight requested values.

\begin{remark}[Why the reciprocal moments are not probability moments]
For $|q|>1$, \eqref{eq:variance} is negative whenever $q>1$.  Therefore
$M_q$ cannot be the MGF of a real probability measure.  This obstruction is
not a defect; it points to the vertical-line and Laplace representation in
\cref{sec:laplace-dual}, where the same numbers become honest complex moments
of a positive real measure.
\end{remark}

\section{The Klein-four sign/inversion calculus}
\label{sec:klein-four}

\subsection{Negative parameters as affine positive laws}

Write $V_j=U_j-1/2$.  The variables $V_j$ are independent and symmetric, so
changing any subset of their signs does not change their joint law.  Since
\begin{equation}
 Y_q-\frac12=(1-q)\sum_{j\ge0}q^jV_j,
 \label{eq:centered-random-series}
\end{equation}
negative geometric ratios are affine images of positive ones.

\begin{theorem}[Negative-$q$ affine conjugacy]
\label{thm:negative-affine}
For $0<a<1$, put
\begin{equation}
 \lambda_a=\frac{1+a}{1-a},
 \qquad
 \beta_a=\frac{a}{1-a}.
 \label{eq:lambda-beta}
\end{equation}
Then
\begin{equation}
 \boxed{
 Y_{-a}\stackrel d=\lambda_aY_a-\beta_a,}
 \qquad
 \boxed{
 Y_{-a}-\frac12\stackrel d=\lambda_a\left(Y_a-\frac12\right).}
 \label{eq:negative-affine}
\end{equation}
Equivalently,
\begin{equation}
 M_{-a}(t)=\e^{-\beta_at}M_a(\lambda_at),
 \qquad
 C_{-a}(t)=C_a(\lambda_at).
 \label{eq:negative-MGF}
\end{equation}
This theorem is \proved{} and \auditnew.
\end{theorem}

\begin{proof}
In \eqref{eq:centered-random-series}, replace $(-1)^jV_j$ by $V_j$ in law:
\[
 Y_{-a}-\frac12
 =(1+a)\sum_{j\ge0}(-a)^jV_j
 \stackrel d=(1+a)\sum_{j\ge0}a^jV_j
 =\lambda_a\left(Y_a-\frac12\right).
\]
Expanding the centered affine map gives the raw shift $-\beta_a$.  Taking
MGFs proves \eqref{eq:negative-MGF}.
\end{proof}

The two examples most relevant here are
\begin{equation}
 Y_{-1/2}\stackrel d=3Y_{1/2}-1,
 \qquad
 Y_{-1/4}\stackrel d=\frac53Y_{1/4}-\frac13.
 \label{eq:negative-special}
\end{equation}
Thus the negative-half law is not a new shape: it is the classical Fabius
density stretched by three and translated to $[-1,2]$.  The negative-quarter
law is the quartic density stretched by $5/3$ and translated to
$[-1/3,4/3]$.

\begin{figure}[H]
\centering
\begin{subfigure}[t]{0.48\textwidth}
\centering
\includegraphics[width=\textwidth]{assets/positive_q_fabius_densities.png}
\caption{The positive laws on $[0,1]$.}
\end{subfigure}\hfill
\begin{subfigure}[t]{0.48\textwidth}
\centering
\includegraphics[width=\textwidth]{assets/negative_q_affine_densities.png}
\caption{Their exact negative-$q$ affine images.}
\end{subfigure}
\caption{Deterministic Fourier inversion of finite characteristic products.
The second panel contains no independent numerical solve: it is obtained from
\cref{eq:negative-affine}.}
\label{fig:density-affine}
\end{figure}

\subsection{Reciprocal inversion}

\begin{theorem}[Reciprocal convolution inversion]
\label{thm:reciprocity}
For every $q$ off the unit circle,
\begin{equation}
 \boxed{M_q(t)M_{1/q}(-t)=1,}
 \qquad
 \boxed{C_q(t)C_{1/q}(t)=1.}
 \label{eq:reciprocity}
\end{equation}
Both equalities hold as analytic germs at $t=0$ and meromorphically wherever
the products are defined.  This theorem is \proved{} and \auditnew.
\end{theorem}

\begin{proof}
Assume first $|q|<1$.  By \eqref{eq:Mq-outside},
\begin{align*}
 M_{1/q}(-t)^{-1}
 &=\prod_{j=1}^{\infty}
 E\!\left((1-q^{-1})q^j(-t)\right)\\
 &=\prod_{j=1}^{\infty}
 E\!\left((1-q)q^{j-1}t\right)
 =M_q(t).
\end{align*}
The case $|q|>1$ follows by interchanging $q$ and $1/q$.  Combining the first
identity with reflection \eqref{eq:M-reflection} gives the centered identity.
\end{proof}

At the coefficient level, \eqref{eq:reciprocity} says
\begin{equation}
 \sum_{k=0}^{n}\binom nk
 m_k(q)(-1)^{n-k}m_{n-k}(1/q)=\delta_{n0}.
 \label{eq:moment-convolution-inverse}
\end{equation}
Thus the reciprocal moment row is the signed binomial-convolution inverse of
the original row.

\subsection{Cayley coordinate and the group action}

Introduce the Cayley coordinate
\begin{equation}
 c(q)=\frac{1+q}{1-q}.
 \label{eq:Cayley}
\end{equation}
A direct calculation gives
\begin{equation}
 c(-q)=c(q)^{-1},
 \qquad
 c(1/q)=-c(q).
 \label{eq:Cayley-action}
\end{equation}
The formal continuation of \cref{thm:negative-affine} can be written entirely
in centered form:
\begin{equation}
 \boxed{C_{-q}(t)=C_q(c(q)t).}
 \label{eq:centered-sign-formal}
\end{equation}
For $q>1$, $c(q)<0$, but $C_q$ is even, so only $|c(q)|$ matters.
Together, \eqref{eq:reciprocity} and \eqref{eq:centered-sign-formal} realize
the Klein four-group by argument dilation and multiplicative inversion.

The two requested orbits therefore reduce to the following explicit formulas:
\begin{align}
 M_{-1/2}(t)&=\e^{-t}M_{1/2}(3t),
 &M_2(t)&=\frac{\e^t}{M_{1/2}(t)},
 &M_{-2}(t)&=\frac{\e^{2t}}{M_{1/2}(3t)},
 \label{eq:dyadic-orbit}\\
 M_{-1/4}(t)&=\e^{-t/3}M_{1/4}(5t/3),
 &M_4(t)&=\frac{\e^t}{M_{1/4}(t)},
 &M_{-4}(t)&=\frac{\e^{4t/3}}{M_{1/4}(5t/3)}.
 \label{eq:quartic-orbit}
\end{align}
Every coefficient identity among the eight requested cases follows from these
six transformations.

\begin{resultbox}{First synthesis principle}
The signs $q\mapsto-q$ and reciprocal scales $q\mapsto1/q$ do not create eight
unrelated special functions.  They create two four-point parameter orbits.
The positive contraction determines the negative probability law by affine
transport and determines both reciprocal germs by convolution inversion.
\end{resultbox}

\section[Spectral q-squared Pochhammer products and generalized Rvachev laws]{Spectral $q^2$-Pochhammer products and generalized Rvachev laws}
\label{sec:spectral}

\subsection{The hyperbolic product of the moment germ}

Euler's product for $\sinh$ gives
\begin{equation}
 S(z)=\frac{\sinh(z/2)}{z/2}
 =\prod_{k=1}^{\infty}\left(1+\frac{z^2}{4\pi^2k^2}\right).
 \label{eq:sinh-product}
\end{equation}
Because $\sum_{j\ge0}(1-q)q^j=1$, the exponential pieces of
\eqref{eq:Mq-inside} combine to $\e^{t/2}$.  Rearranging the remaining
normally convergent product yields the spectral factorization.

\begin{theorem}[Spectral Pochhammer factorization]
\label{thm:spectral-Pochhammer}
For $|q|<1$,
\begin{equation}
 \boxed{
 M_q(t)=\e^{t/2}
 \prod_{k=1}^{\infty}
 \qpoch{-\dfrac{(1-q)^2t^2}{4\pi^2k^2}}{q^2}{\infty}.}
 \label{eq:M-Pochhammer-inside}
\end{equation}
For $|q|>1$,
\begin{equation}
 \boxed{
 M_q(t)=\e^{t/2}
 \prod_{k=1}^{\infty}
 \qpoch{-\dfrac{(1-q^{-1})^2t^2}{4\pi^2k^2}}{q^{-2}}{\infty}^{-1}.}
 \label{eq:M-Pochhammer-outside}
\end{equation}
The theorem is \proved.
\end{theorem}

\begin{proof}
Insert \eqref{eq:sinh-product} into
$E((1-q)q^jt)=\e^{(1-q)q^jt/2}S((1-q)q^jt)$.  For fixed $k$, the product over
$j$ is
\[
 \prod_{j\ge0}\left(1+
 \frac{(1-q)^2q^{2j}t^2}{4\pi^2k^2}\right)
 =\qpoch{-\dfrac{(1-q)^2t^2}{4\pi^2k^2}}{q^2}{\infty}.
\]
This proves \eqref{eq:M-Pochhammer-inside}.  In the reciprocal product, set
$r=q^{-1}$.  The scale $(1-q)q^{-j}=-(1-r)r^{j-1}$; its square and the inverse
product give \eqref{eq:M-Pochhammer-outside}.  The exponential normalization
is again $\e^{t/2}$.
\end{proof}

This theorem resolves an apparent ambiguity in the requested parameters.  The
geometric ratio is $q$, but the spectral Pochhammer nome is $q^2$.  Therefore
\[
 q=\pm\frac12 \quad\Longrightarrow\quad \text{spectral nome }\frac14,
 \qquad
 q=\pm\frac14 \quad\Longrightarrow\quad \text{spectral nome }\frac1{16}.
\]
Their reciprocal partners $\pm2$ and $\pm4$ use the same contracting nomes but
with inverse factors.  Inversion is literally a zero--pole exchange.

\subsection{The centered Rvachev random series}

Let
\begin{equation}
 R_q=\sum_{j=0}^{\infty}q^jV_j,
 \qquad V_j=U_j-\frac12,
 \qquad |q|<1.
 \label{eq:Rq}
\end{equation}
Then
\begin{equation}
 Y_q=\frac12+(1-q)R_q.
 \label{eq:Y-R}
\end{equation}
Under the Fourier convention
$\widehat\mu(z)=\int\e^{2\pi\ii zx}\,\dd\mu(x)$, its characteristic function is
\begin{equation}
 \Phi_q(z)=\E\e^{2\pi\ii zR_q}
 =\prod_{j=0}^{\infty}\sinc(q^jz),
 \qquad
 \sinc z=\frac{\sin\pi z}{\pi z}.
 \label{eq:Phi-q}
\end{equation}
At $q=1/2$, $R_q$ is the random variable whose density is Rvachev's
up-function in the standard support normalization.  Since the $V_j$ are
symmetric,
\begin{equation}
 R_{-a}\stackrel d=R_a,
 \qquad 0<a<1.
 \label{eq:R-sign}
\end{equation}
Thus the sign dependence of $Y_q$ is entirely caused by the outer affine
normalization $1-q$.

The Fourier product itself has the Pochhammer form
\begin{equation}
 \boxed{
 \Phi_q(z)=\prod_{k=1}^{\infty}
 \qpoch{z^2/k^2}{q^2}{\infty}.}
 \label{eq:Phi-Pochhammer}
\end{equation}
Its Mahler equation is
\begin{equation}
 \Phi_q(z)=\sinc(z)\Phi_q(qz).
 \label{eq:Phi-Mahler}
\end{equation}

\subsection{Integer reciprocal scales and zero multiplicities}

The classical dyadic transform has zero multiplicity $1+\nu_2(n)$ at every
nonzero integer $n$.  The same argument works for every integer reciprocal
scale.

\begin{proposition}[Base-$b$ zero divisor]
\label{prop:base-b-zero}
Let $q=1/b$ with integer $b\ge2$.  At a nonzero integer $n$,
\begin{equation}
 \ord_{z=n}\Phi_{1/b}(z)=1+\nu_b(n),
 \label{eq:base-b-zero-order}
\end{equation}
where $\nu_b(n)=\max\{j:b^j\mid n\}$.  This is \proved.
\end{proposition}

\begin{proof}
The factor $\sinc(z/b^j)$ has a simple zero at $z=n$ exactly when
$n/b^j$ is a nonzero integer, that is, when $0\le j\le\nu_b(n)$.  All other
factors are nonzero there.
\end{proof}

For $b=4$,
\begin{equation}
 \ord_{z=n}\Phi_{1/4}(z)
 =1+\floor{\frac{\nu_2(n)}2}.
 \label{eq:quartic-zero-order}
\end{equation}
The quartic transform therefore sees pairs of binary valuation levels at once.
This is the spectral counterpart of the even/odd scale decimation in the next
section.

\section[Scale multisection and q-integer convolution]{Scale multisection and $q$-integer convolution}
\label{sec:multisection}

\subsection[Decomposing scales modulo m]{Decomposing scales modulo $m$}

Define the $q$-integer
\begin{equation}
 [m]_q=1+q+\cdots+q^{m-1}=\frac{1-q^m}{1-q}.
 \label{eq:q-integer}
\end{equation}
Separate the coordinate index in \eqref{eq:Yq} as $j=mk+r$.

\begin{theorem}[Geometric multisection]
\label{thm:geometric-multisection}
Let $m\ge2$, $|q|\ne1$, and $[m]_q\ne0$.  As analytic germs,
\begin{equation}
 \boxed{
 M_q(t)=\prod_{r=0}^{m-1}
 M_{q^m}\!\left(\frac{q^r}{[m]_q}t\right).}
 \label{eq:M-multisection}
\end{equation}
For $|q|<1$, this is the distributional identity
\begin{equation}
 \boxed{
 Y_q\stackrel d=
 \sum_{r=0}^{m-1}\frac{q^r}{[m]_q}Y_{q^m}^{(r)},}
 \label{eq:Y-multisection}
\end{equation}
where the variables on the right are independent.  The theorem is \proved{}
and \auditnew{} in this normalized form.
\end{theorem}

\begin{proof}
For $|q|<1$,
\begin{align*}
 Y_q
 &=(1-q)\sum_{r=0}^{m-1}\sum_{k\ge0}q^{mk+r}U_{mk+r}\\
 &=\sum_{r=0}^{m-1}
 \frac{(1-q)q^r}{1-q^m}
 \left((1-q^m)\sum_{k\ge0}(q^m)^kU_{mk+r}\right),
\end{align*}
which is \eqref{eq:Y-multisection}.  Independence gives
\eqref{eq:M-multisection}.  The two sides are normalized formal solutions and
the identity extends to the reciprocal regime by the product definition, or
directly by multisection of the contracting reciprocal product.
\end{proof}

For $m=2$,
\begin{equation}
 Y_q\stackrel d=
 \frac1{1+q}Y_{q^2}^{(0)}+
 \frac q{1+q}Y_{q^2}^{(1)}.
 \label{eq:Y-decimation-two}
\end{equation}
The requested specializations include
\begin{align}
 Y_{1/2}&\stackrel d=\frac23Y_{1/4}^{(0)}+\frac13Y_{1/4}^{(1)},
 \label{eq:Fabius-quartic-convolution}\\
 Y_{-1/2}&\stackrel d=2Y_{1/4}^{(0)}-Y_{1/4}^{(1)},
 \label{eq:negative-half-quartic}\\
 M_2(t)&=M_4(t/3)M_4(2t/3).
 \label{eq:reciprocal-decimation}
\end{align}
The first identity decomposes the classical Fabius distribution into two
independent quartic distributions at unequal scales.  The second is a signed
version of the same decomposition and agrees with
$Y_{-1/2}\stackrel d=3Y_{1/2}-1$.

\subsection{Rvachev multisection}

The unnormalized centered series has the cleaner identity
\begin{equation}
 \boxed{
 R_q\stackrel d=\sum_{r=0}^{m-1}q^rR_{q^m}^{(r)},}
 \qquad
 \boxed{
 \Phi_q(z)=\prod_{r=0}^{m-1}\Phi_{q^m}(q^rz).}
 \label{eq:R-Phi-multisection}
\end{equation}
At $q=1/2$, $m=2$, the dyadic up-density is therefore a convolution of two
quartic up-densities, one compressed by two.  In density notation,
\begin{equation}
 u_{1/2}=u_{1/4}*\left(2u_{1/4}(2\,\cdot)\right).
 \label{eq:up-quartic-convolution}
\end{equation}
At the product level, all these identities are instances of the elementary
Pochhammer multisection
\begin{equation}
 (a;q)_\infty=\prod_{r=0}^{m-1}(aq^r;q^m)_\infty.
 \label{eq:Pochhammer-multisection}
\end{equation}

\begin{resultbox}{Second synthesis principle}
The classical dyadic object is not isolated from the quartic object.  It is an
exact two-component weighted independent-sum decomposition at the random-variable
level, a convolution
at the density level, a factorization at the Fourier level, and a multisection
identity at the Pochhammer level.  Negative and reciprocal parameters preserve
the same diagram, with signed coefficients or inverse factors as appropriate.
\end{resultbox}

\section{Reciprocal Laplace duality and vertical moment geometry}
\label{sec:laplace-dual}

\subsection{From inverse hyperbolic products to Laplace laws}

The negative formal variance for $q>1$ suggests that the reciprocal germ should
not be sought on the real line.  The centered identity
$C_q=1/C_{1/q}$ turns that obstruction into a positive characteristic
function.

Let $r=1/q$, so $|r|<1$, and define
\begin{equation}
 b_{j,k}(r)=\frac{|1-r|\,|r|^j}{2\pi k},
 \qquad j\ge0,\ k\ge1.
 \label{eq:b-jk}
\end{equation}
Let $Z_{j,k}$ be independent standard symmetric Laplace variables with density
$\frac12\e^{-|x|}$ and characteristic function
$(1+t^2)^{-1}$.

\begin{theorem}[Positive Laplace representation of the reciprocal germ]
\label{thm:laplace-dual}
For every real $q$ with $|q|>1$, the series
\begin{equation}
 L_{1/q}=\sum_{j=0}^{\infty}\sum_{k=1}^{\infty}
 b_{j,k}(1/q)Z_{j,k}
 \label{eq:Lr}
\end{equation}
converges in $L^2$ and almost surely, and
\begin{equation}
 \boxed{C_q(t)=\E\e^{\ii tL_{1/q}}.}
 \label{eq:C-Laplace}
\end{equation}
Consequently
\begin{equation}
 \boxed{m_n(q)=\E\left(\frac12+\ii L_{1/q}\right)^n.}
 \label{eq:vertical-moments}
\end{equation}
This theorem is \proved{} and \auditnew.
\end{theorem}

\begin{proof}
By \eqref{eq:M-Pochhammer-inside} and reciprocity,
\begin{align*}
 C_q(t)
 &=\frac1{C_r(t)}\\
 &=\prod_{j\ge0}\prod_{k\ge1}
 \left(1+\frac{(1-r)^2r^{2j}t^2}{4\pi^2k^2}\right)^{-1}\\
 &=\prod_{j,k}\left(1+b_{j,k}(r)^2t^2\right)^{-1}.
\end{align*}
Moreover,
\[
 \sum_{j,k}b_{j,k}(r)^2
 =\frac{(1-r)^2}{4\pi^2(1-r^2)}\zeta(2)<\infty.
\]
Since a standard Laplace variable has variance $2$, Kolmogorov's convergence
criterion gives $L^2$ and almost-sure convergence of \eqref{eq:Lr}.  The
characteristic functions of the finite partial sums converge to the displayed
product, proving \eqref{eq:C-Laplace}.  Finally,
$M_q(t)=\e^{t/2}C_q(t)=\E\exp(t(1/2+\ii L_r))$ near the origin, and coefficient
comparison proves \eqref{eq:vertical-moments}.
\end{proof}

The theorem replaces the phrase ``signed probability'' by a precise geometric
object: a positive measure on the vertical line
\begin{equation}
 \left\{\frac12+\ii x:x\in\R\right\}.
 \label{eq:vertical-line}
\end{equation}
The reciprocal moments are ordinary polynomial moments of that measure.

\subsection{The zeta bridge in the Laplace cumulants}

For a Laplace variable of scale $b$, the even cumulants are
$(2m)!b^{2m}/m$.  Hence
\begin{equation}
 \kappa_{2m}(L_r)=
 \frac{(2m)!}{m}
 \frac{(1-r)^{2m}}{(2\pi)^{2m}(1-r^{2m})}\zeta(2m).
 \label{eq:Laplace-cumulants}
\end{equation}
Using
\begin{equation}
 B_{2m}=(-1)^{m+1}\frac{2(2m)!}{(2\pi)^{2m}}\zeta(2m)
 \label{eq:Bernoulli-zeta}
\end{equation}
and $q=1/r$, one recovers
\begin{equation}
 \kappa_{2m}(q)=(-1)^m\kappa_{2m}(L_r),
 \label{eq:vertical-cumulants}
\end{equation}
exactly as required by the factor $\ii^{2m}$ in
\eqref{eq:vertical-moments}.  This is a direct three-way bridge among
Bernoulli numbers, the Riemann zeta values, and the reciprocal $q$-Fabius
germ.

\subsection{Exact Hankel signature}

Let
\begin{equation}
 \Delta_n(q)=\det\left(m_{i+j}(q)\right)_{0\le i,j<n}.
 \label{eq:Hankel-q}
\end{equation}

\begin{theorem}[Hankel signature]
\label{thm:Hankel-signature}
For every real $|q|>1$ and $n\ge1$,
\begin{equation}
 \boxed{
 \sgn\Delta_n(q)=(-1)^{\binom n2}.}
 \label{eq:Hankel-signature}
\end{equation}
In particular, no Hankel determinant vanishes.  The theorem is \proved{} and
\auditnew.
\end{theorem}

\begin{proof}
Let $H_n(L_r)$ be the positive Hankel determinant of the real moments of
$L_r$.  It is strictly positive because the Laplace convolution has infinite
support, so the monomials $1,x,\ldots,x^{n-1}$ are linearly independent in
$L^2$ of its law.  Translation of a moment sequence leaves its Hankel
determinant unchanged: the change-of-basis matrix from $x^k$ to
$(x+c)^k$ is unit lower triangular.  Scaling the variable by $a$ multiplies an
$n$th Hankel determinant by $a^{n(n-1)}$.  Applying these facts to
$1/2+\ii L_r$ gives
\[
 \Delta_n(q)=\ii^{n(n-1)}H_n(L_r)
 =(-1)^{\binom n2}H_n(L_r).
\]
\end{proof}

The exact determinants through order nine for $q=\pm2,\pm4$ are included in
\path{data/reciprocal_hankel_signatures.csv}.  Their sizes vary enormously,
but their signs follow \eqref{eq:Hankel-signature} without exception.

\subsection{Orthogonal polynomials on a vertical line}

Let $Q_n^{(r)}$ be the monic orthogonal polynomial of degree $n$ for the real
law of $L_r$.  Symmetry gives a recurrence
\begin{equation}
 Q_{n+1}^{(r)}(x)=xQ_n^{(r)}(x)-\beta_n^{(r)}Q_{n-1}^{(r)}(x),
 \qquad \beta_n^{(r)}>0.
 \label{eq:Q-recurrence}
\end{equation}
Define
\begin{equation}
 P_n^{(q)}(z)=\ii^nQ_n^{(1/q)}\!\left(\frac{z-1/2}{\ii}\right).
 \label{eq:vertical-OP}
\end{equation}
Then $P_n^{(q)}$ is monic and orthogonal for the reciprocal moment functional.
Its recurrence is
\begin{equation}
 P_{n+1}^{(q)}(z)
 =(z-1/2)P_n^{(q)}(z)+\beta_n^{(1/q)}P_{n-1}^{(q)}(z).
 \label{eq:P-recurrence}
\end{equation}
In the standard monic form $P_{n+1}=(z-a_n)P_n-b_nP_{n-1}$, the subdiagonal
coefficient is $b_n=-\beta_n<0$.  Every zero has the form
\begin{equation}
 z=\frac12+\ii x,
 \qquad x\in\R.
 \label{eq:vertical-zeros}
\end{equation}
This supplies a concrete spectral problem for the reciprocal germs: determine
the recurrence coefficients $\beta_n^{(r)}$, their asymptotics, and their
relation to the geometric double lattice $b_{j,k}(r)$.

\begin{resultbox}{Third synthesis principle}
Reciprocal $q$ does not destroy positivity; it rotates it.  The centered
reciprocal germ is positive definite on the Fourier axis, the raw moments live
on $\Re z=1/2$, the Hankel form has a deterministic alternating signature,
and its orthogonal zeros align on the same vertical line.
\end{resultbox}

\section[q-Fabius--Bernoulli Appell polynomials]{$q$-Fabius--Bernoulli Appell polynomials}
\label{sec:appell}

\subsection{Inverse geometric averaging}

For $|r|<1$, define the translation-averaging operator on polynomials by
\begin{equation}
 (\mathcal A_r f)(x)=\E f(x+Y_r).
 \label{eq:averaging-operator}
\end{equation}
On formal exponentials, $\mathcal A_r$ is the multiplier $M_r(D)$:
\[
 \mathcal A_r\e^{xt}=M_r(t)\e^{xt}.
\]
Reciprocity therefore identifies its inverse without solving a triangular
system coefficient by coefficient:
\begin{equation}
 \boxed{
 \mathcal A_r^{-1}=M_{1/r}(-D).}
 \label{eq:averaging-inverse}
\end{equation}
The reciprocal $q$-Fabius germ is precisely the deconvolution symbol for
geometric-uniform averaging.

\begin{definition}[$q$-Fabius--Bernoulli polynomials]
For $|r|<1$, define the Appell sequence $\Br_n^{(r)}(x)$ by
\begin{equation}
 \boxed{
 \sum_{n=0}^{\infty}\Br_n^{(r)}(x)\frac{t^n}{n!}
 =\frac{\e^{xt}}{M_r(t)}.}
 \label{eq:Fabius-Bernoulli-gen}
\end{equation}
\end{definition}

\begin{theorem}[Exact deconvolution]
\label{thm:Appell-deconvolution}
For every polynomial degree $n$,
\begin{equation}
 \boxed{
 \E\Br_n^{(r)}(x+Y_r)=x^n.}
 \label{eq:Appell-deconvolution}
\end{equation}
The sequence is Appell:
\begin{equation}
 \frac{\dd}{\dd x}\Br_n^{(r)}(x)=n\Br_{n-1}^{(r)}(x).
 \label{eq:Appell-derivative}
\end{equation}
These statements are \proved{} and \auditnew.
\end{theorem}

\begin{proof}
Taking the expectation of the generating function after replacing $x$ by
$x+Y_r$ multiplies its right-hand side by $M_r(t)$, leaving $\e^{xt}$.
Coefficient comparison proves \eqref{eq:Appell-deconvolution}.  Differentiating
\eqref{eq:Fabius-Bernoulli-gen} with respect to $x$ proves
\eqref{eq:Appell-derivative}.
\end{proof}

\subsection{Bernoulli and concentration limits}

At $r=0$, $Y_0=U_0$ and $M_0(t)=E(t)$.  Therefore
\begin{equation}
 \frac{\e^{xt}}{M_0(t)}=\frac{t\e^{xt}}{\e^t-1},
 \qquad
 \boxed{\Br_n^{(0)}(x)=B_n(x),}
 \label{eq:Bernoulli-boundary}
\end{equation}
where $B_n(x)$ is the ordinary Bernoulli polynomial with
$B_1(x)=x-1/2$.  At the opposite probability boundary,
$\Var(Y_r)\to0$ as $r\to1^-$ and $M_r(t)\to\e^{t/2}$ locally uniformly.
Hence
\begin{equation}
 \boxed{
 \Br_n^{(r)}(x)\longrightarrow(x-1/2)^n
 \quad(r\to1^-).}
 \label{eq:Appell-concentration}
\end{equation}
The family continuously interpolates Bernoulli deconvolution of one uniform
average and trivial deconvolution of a point mass at $1/2$.

Reciprocity also explains the boundary from outside the unit disk.  As
$|q|\to\infty$, $1/q\to0$ and
\begin{equation}
 M_q(t)=\frac{\e^t}{M_{1/q}(t)}
 \longrightarrow \frac{t\e^t}{\e^t-1}
 =\sum_{n\ge0}B_n(1)\frac{t^n}{n!}.
 \label{eq:reciprocal-Bernoulli-limit}
\end{equation}
Thus Bernoulli polynomials sit at both ends of the sign/reciprocal diagram:
as the inverse of ordinary uniform averaging and as the limiting reciprocal
moment germ.

\subsection{A Bernoulli--Mahler refinement equation}

For $0<|r|<1$, use the Mahler equation $M_r(t)=E((1-r)t)M_r(rt)$ in
\eqref{eq:Fabius-Bernoulli-gen}:
\begin{align}
 \frac{\e^{xt}}{M_r(t)}
 &=\frac{(1-r)t}{\e^{(1-r)t}-1}
   \frac{\e^{xt}}{M_r(rt)}\nonumber\\
 &=\left(\sum_{k\ge0}B_k\frac{((1-r)t)^k}{k!}\right)
   \left(\sum_{j\ge0}\Br_j^{(r)}(x/r)\frac{(rt)^j}{j!}\right).
 \label{eq:Appell-Mahler-gen}
\end{align}
Here $B_1=-1/2$ is the standard Bernoulli-number convention.  Therefore
\begin{equation}
 \boxed{
 \Br_n^{(r)}(x)=
 \sum_{k=0}^{n}\binom nk
 B_k(1-r)^kr^{n-k}\Br_{n-k}^{(r)}(x/r).}
 \label{eq:Appell-Mahler}
\end{equation}
This relation is a Bernoulli-polynomial analogue of the Fabius Mahler equation.
It combines a local uniform deconvolution at scale $1-r$ with the same Appell
family at the contracted argument $x/r$.

\subsection{Bell-polynomial form and first examples}

Taking logarithms in \eqref{eq:Fabius-Bernoulli-gen} gives
\begin{equation}
 \log\frac{\e^{xt}}{M_r(t)}
 =(x-1/2)t-
 \sum_{m\ge1}\kappa_{2m}(r)\frac{t^{2m}}{(2m)!}.
 \label{eq:Appell-log}
\end{equation}
Consequently
\begin{equation}
 \boxed{
 \Br_n^{(r)}(x)=
 \Bell_n\left(x-\frac12,-\kappa_2(r),0,-\kappa_4(r),0,\ldots\right).}
 \label{eq:Appell-Bell}
\end{equation}
Writing $z=x-1/2$ and
$v_r=\kappa_2(r)=(1-r)/(12(1+r))$, one obtains
\begin{align}
 \Br_0^{(r)}(x)&=1,
 &\Br_1^{(r)}(x)&=z,\nonumber\\
 \Br_2^{(r)}(x)&=z^2-v_r,
 &\Br_3^{(r)}(x)&=z^3-3v_rz,\nonumber\\
 \Br_4^{(r)}(x)&=z^4-6v_rz^2+3v_r^2-\kappa_4(r).
 \label{eq:Appell-first}
\end{align}
The symmetry
\begin{equation}
 \Br_n^{(r)}(1-x)=(-1)^n\Br_n^{(r)}(x)
 \label{eq:Appell-reflection}
\end{equation}
follows immediately from centered evenness.  Their arithmetic, root geometry,
and possible analogues of the von Staudt--Clausen and multiplication theorems
form a natural new research program.

\section{A Fabius--Pochhammer moment polynomial}
\label{sec:moment-polynomial}

\subsection{Definition and polynomial recurrence}

The rational moments $a_n(q)=m_n(q)/n!$ have poles at roots of unity because of
the factors $1-q^j$ in the recurrence.  Multiplying by the finite
$q$-Pochhammer symbol and removing the obvious $(1-q)^n$ scale produces a
polynomial.

\begin{definition}
For $n\ge0$, define
\begin{equation}
 \boxed{
 \mathcal P_n(q)=
 \frac{(q;q)_n}{(1-q)^n}\frac{m_n(q)}{n!}.}
 \label{eq:Pn-def}
\end{equation}
At $q=1$ the right-hand expression is interpreted by polynomial continuation.
\end{definition}

\begin{theorem}[Polynomial recurrence]
\label{thm:Pn-recurrence}
The functions $\mathcal P_n$ are polynomials in $\Q[q]$, with
$\deg\mathcal P_n\le\binom n2$, and
\begin{equation}
 \mathcal P_0(q)=1,
 \label{eq:P0}
\end{equation}
\begin{equation}
 \boxed{
 \mathcal P_n(q)=
 \sum_{k=0}^{n-1}
 \frac{q^k}{(n-k+1)!}
 \frac{(q;q)_{n-1}}{(q;q)_k}\mathcal P_k(q).}
 \label{eq:Pn-recurrence}
\end{equation}
Moreover,
\begin{equation}
 \mathcal P_n(0)=\frac1{(n+1)!},
 \qquad
 [q^{\binom n2}]\mathcal P_n(q)=\frac{B_n(1)}{n!}.
 \label{eq:Pn-boundaries}
\end{equation}
The theorem is \proved{} and \auditnew.
\end{theorem}

\begin{proof}
Multiply \eqref{eq:an-recurrence} by
$(q;q)_{n-1}/(1-q)^n$ and substitute
$a_k=(1-q)^k\mathcal P_k/(q;q)_k$.  This gives
\eqref{eq:Pn-recurrence}.  The quotient
$(q;q)_{n-1}/(q;q)_k=\prod_{j=k+1}^{n-1}(1-q^j)$ is polynomial.  If
$\deg\mathcal P_k\le\binom k2$, the $k$th summand has degree at most
\[
 k+\sum_{j=k+1}^{n-1}j+\binom k2=\binom n2,
\]
so induction proves polynomiality and the degree bound.

At $q=0$, only the $k=0$ summand survives, giving
$\mathcal P_n(0)=1/(n+1)!$.  For the leading coefficient, let $q\to\infty$.
By \eqref{eq:reciprocal-Bernoulli-limit},
$m_n(q)/n!\to B_n(1)/n!$, while
\[
 \frac{(q;q)_n}{(1-q)^n}
 \sim q^{\binom n2}.
\]
The coefficient assertion follows.
\end{proof}

The first polynomials are
\begin{align}
 \mathcal P_0(q)&=1,
 &\mathcal P_1(q)&=\frac12,\nonumber\\
 \mathcal P_2(q)&=\frac{q+2}{12},
 &\mathcal P_3(q)&=\frac{q^2+q+1}{24},
 \label{eq:Pn-first-a}\\
 \mathcal P_4(q)&=
 -\frac{(q^2+q+1)(q^4-2q^3-5q^2-3q-6)}{720}.
 \label{eq:Pn-first-b}
\end{align}
The vanishing of $B_n(1)$ for odd $n>1$ explains the strict degree drop
$\deg\mathcal P_n<\binom n2$ in odd degrees.

\subsection{A cyclotomic divisibility pattern}

The exact computations through degree $16$ reveal a stronger pattern:
\begin{equation}
 \prod_{\substack{3\le d\le n\\d\ \mathrm{odd}}}[d]_q
 \quad\text{appears to divide}\quad \mathcal P_n(q),
 \qquad
 [d]_q=1+q+\cdots+q^{d-1}.
 \label{eq:Pn-divisibility-observed}
\end{equation}
It is visible already in \eqref{eq:Pn-first-a}--\eqref{eq:Pn-first-b}, and
continues in the computed factorizations of $\mathcal P_5,\ldots,\mathcal P_{16}$.
The exact polynomial, proposed divisor, and quotient data are recorded in
\path{data/moment_polynomial_divisibility.csv}.
The odd-cumulant vanishing suggests the mechanism, but common cyclotomic
factors between $1-q^{2m}$ and $1-q^d$ make a direct denominator-counting
argument insufficient.

\begin{conjecture}[Odd $q$-integer divisor]
\label{conj:Pn-divisibility}
For every $n\ge3$,
\begin{equation}
 \boxed{
 \prod_{\substack{3\le d\le n\\d\ \mathrm{odd}}}[d]_q
 \mid \mathcal P_n(q)
 \quad\text{in }\Q[q].}
 \label{eq:Pn-divisibility}
\end{equation}
This is \openstatus.  It has been checked symbolically for $n\le16$.
\end{conjecture}

A proof may be accessible through a root-of-unity analysis of the centered
Bell expansion, or through a cyclic version of the two-nome partition function
in \cref{sec:two-nome}.  The multiplicities of repeated cyclotomic factors in
the product make the conjecture substantially stronger than separate
vanishing at primitive odd roots.

\subsection{Endpoint interpretation}

The main reason for introducing $\mathcal P_n$ is that the exact endpoint
formula of \cref{sec:endpoints} becomes
\begin{equation}
 \boxed{
 G_q(q^n)=\frac{q^{\binom{n+1}{2}}}{(q;q)_n}\mathcal P_n(q)
 \qquad(0<q\le1/2).}
 \label{eq:endpoint-Pn-preview}
\end{equation}
Thus the finite $q$-Pochhammer denominator at the inverse-geometric endpoint is
not an imported notation: it is the exact pole-clearing factor of the moment
recurrence.  At $q=1/2$ it becomes the familiar product of Mersenne factors;
at $q=1/4$ it becomes a product of $4^j-1$ factors.

\section{A two-nome Pochhammer--Prouhet partition function}
\label{sec:two-nome}

\subsection{Definition on the binary cube}

For a binary vector $\epsilon=(\epsilon_0,\ldots,\epsilon_{N-1})\in\{0,1\}^N$,
define its Hamming weight, position energy, and geometric projection by
\begin{equation}
 |\epsilon|=\sum_{j=0}^{N-1}\epsilon_j,
 \qquad
 \Omega(\epsilon)=\sum_{j=0}^{N-1}j\epsilon_j,
 \qquad
 A_{\sigma,N}(\epsilon)=\sum_{j=0}^{N-1}\epsilon_j\sigma^j.
 \label{eq:binary-statistics}
\end{equation}
The usual Thue--Morse sign on the cube is $(-1)^{|\epsilon|}$.

\begin{definition}[Two-nome binary-cube partition function]
\label{def:two-nome}
For parameters $u,t,\rho,\sigma$, set
\begin{equation}
 \boxed{
 \Zpart_N(u,t;\rho,\sigma)=
 \prod_{j=0}^{N-1}\left(1-u\rho^j\e^{t\sigma^j}\right).}
 \label{eq:Z-master}
\end{equation}
\end{definition}

Expanding one factor at a time gives the exact cube sum
\begin{equation}
 \boxed{
 \Zpart_N(u,t;\rho,\sigma)
 =\sum_{\epsilon\in\{0,1\}^N}
 (-u)^{|\epsilon|}\rho^{\Omega(\epsilon)}
 \e^{tA_{\sigma,N}(\epsilon)}.}
 \label{eq:Z-cube}
\end{equation}
The four variables have distinct jobs:
$u$ marks Hamming weight, $\rho$ marks position energy, $t$ marks projected
moments, and $\sigma$ chooses the geometric projection.  This separation is
the cleanest way to connect Gaussian coefficients to Thue--Morse Prouhet sums
without identifying their nomes by fiat.

\subsection{Two boundary faces}

At $t=0$,
\begin{equation}
 \Zpart_N(u,0;\rho,\sigma)=(u;\rho)_N.
 \label{eq:Z-Pochhammer-face}
\end{equation}
The finite $q$-binomial theorem then gives
\begin{equation}
 (u;\rho)_N
 =\sum_{k=0}^{N}(-u)^k\rho^{\binom k2}
 \qbinom{N}{k}{\rho}.
 \label{eq:finite-q-binomial}
\end{equation}
At the opposite face $(u,\rho)=(1,1)$,
\begin{equation}
 \Zpart_N(1,t;1,\sigma)
 =\prod_{j=0}^{N-1}\left(1-\e^{t\sigma^j}\right)
 =\sum_{\epsilon}(-1)^{|\epsilon|}
 \e^{tA_{\sigma,N}(\epsilon)}.
 \label{eq:Z-Prouhet-face}
\end{equation}
This is the exponential generating function of the geometric Prouhet moments.
Thus the $q$-Pochhammer and Thue--Morse products are not merely analogous:
they are perpendicular faces of one four-parameter finite product.

\begin{resultbox}{Master binary-cube object}
The finite Gaussian coefficient records how many selected digit positions
carry a given total energy.  The Prouhet generator records how the same subset
projects to a geometric number.  The partition function $\Zpart_N$ retains both
statistics simultaneously and permits independent nomes $\rho$ and $\sigma$.
\end{resultbox}

\subsection{Simultaneous multisection}

Splitting the index $j=mk+r$ gives a multisection in both nomes at once.

\begin{theorem}[Two-nome multisection]
\label{thm:Z-multisection}
For integers $m,N\ge1$,
\begin{equation}
 \boxed{
 \Zpart_{mN}(u,t;\rho,\sigma)
 =\prod_{r=0}^{m-1}
 \Zpart_N\!\left(u\rho^r,t\sigma^r;\rho^m,\sigma^m\right).}
 \label{eq:Z-multisection}
\end{equation}
The identity is \proved{} and \auditnew.
\end{theorem}

\begin{proof}
The factor with index $j=mk+r$ on the left is
\[
 1-u\rho^{mk+r}\e^{t\sigma^{mk+r}}
 =1-(u\rho^r)(\rho^m)^k
   \e^{(t\sigma^r)(\sigma^m)^k},
\]
which is the $k$th factor of the $r$th product on the right.
\end{proof}

At $t=0$, \eqref{eq:Z-multisection} is finite Pochhammer multisection.  At
$(u,\rho)=(1,1)$, it is the derivative-comb convolution law developed in
\cref{sec:comb}.  At $N\to\infty$ in the normalized Prouhet jet, it becomes the
geometric-uniform multisection theorem \ref{thm:geometric-multisection}.

\subsection{Mixed Gaussian--Prouhet moments}

Define
\begin{equation}
 \mathcal M_{N,k,m}(\rho,\sigma)=
 \sum_{|\epsilon|=k}
 \rho^{\Omega(\epsilon)}A_{\sigma,N}(\epsilon)^m.
 \label{eq:mixed-moments}
\end{equation}
From \eqref{eq:Z-cube},
\begin{equation}
 \boxed{
 \mathcal M_{N,k,m}(\rho,\sigma)
 =(-1)^k m!\,[u^kt^m]\Zpart_N(u,t;\rho,\sigma).}
 \label{eq:mixed-coefficient}
\end{equation}
This coefficient simultaneously refines a Gaussian row by geometric power
moments.  For $m=0$ it is the position-energy layer
$\rho^{\binom k2}\qbinom{N}{k}{\rho}$.  Summing over $k$ with sign recovers the
Prouhet moment in \cref{sec:q-prouhet}.

\subsection{Root-of-unity cancellation beyond classical Prouhet}

The order of $\Zpart_N$ at $t=0$ can be read factor by factor.

\begin{proposition}[Cancellation order]
\label{prop:Z-order}
Assume $\sigma\ne0$.  Then
\begin{equation}
 \ord_{t=0}\Zpart_N(u,t;\rho,\sigma)
 =\#\left\{0\le j<N:u\rho^j=1\right\}.
 \label{eq:Z-order}
\end{equation}
Consequently the corresponding weighted mixed moments vanish below this
order.  The proposition is \proved.
\end{proposition}

\begin{proof}
At $t=0$, the $j$th factor is $1-u\rho^j$.  If it is nonzero, it contributes
no zero.  If $u\rho^j=1$, the factor is
$1-\e^{t\sigma^j}=-\sigma^jt+O(t^2)$ and has one simple zero.  Orders add under
multiplication.
\end{proof}

Classical Prouhet cancellation is the maximal case $u=\rho=1$, in which all
$N$ factors vanish.  At $u=1$, $\rho=-1$, only even positions contribute
zeros, so the order is $\ceil{N/2}$.  More generally, root-of-unity choices of
$\rho$ select arithmetic progressions of digit positions.  This gives a
finite, exactly controlled family of ``partial Prouhet'' cancellations that is
absent if the Gaussian and geometric nomes are collapsed too early.

\subsection{The Fabius MGF as a renormalized Prouhet jet}

Set $(u,\rho,\sigma)=(1,1,q)$ and replace $t$ by $(1-q)t$.  Since
$1-\e^z=-zE(z)$,
\begin{align}
 \Zpart_N(1,(1-q)t;1,q)
 &=(-1)^N(1-q)^Nq^{\binom N2}t^N
   \prod_{j=0}^{N-1}E((1-q)q^jt).
 \label{eq:Z-renormalized-finite}
\end{align}
Therefore, for $0<|q|<1$,
\begin{equation}
 \boxed{
 M_q(t)=\lim_{N\to\infty}
 \frac{\Zpart_N(1,(1-q)t;1,q)}
 {(-1)^N(1-q)^Nq^{\binom N2}t^N}.}
 \label{eq:M-renormalized-Prouhet}
\end{equation}
The Fabius MGF is the normalized residual jet left after removing the maximal
order-$N$ Prouhet zero.  This identity makes the transition from a finite
Thue--Morse cancellation block to the smooth infinite Fabius distribution
completely explicit.

\section{Signed position-energy layers and negative Gaussian coefficients}
\label{sec:position-energy}

\subsection{The Gaussian layer formula}

At fixed Hamming weight $k$, define
\begin{equation}
 H_{N,k}(\rho)=
 \sum_{|\epsilon|=k}\rho^{\Omega(\epsilon)}.
 \label{eq:H-Nk}
\end{equation}
Extracting $u^k$ in \eqref{eq:finite-q-binomial} gives
\begin{equation}
 \boxed{
 H_{N,k}(\rho)=
 \rho^{\binom k2}\qbinom{N}{k}{\rho}.}
 \label{eq:H-Gaussian}
\end{equation}
This is the principal specialization of the Schur function of a rectangular
partition, but its binary-cube meaning is more useful here: it is the exact
partition function of $k$ selected bit positions with linear energy $j$.

Complementing selected positions or reversing them gives the reciprocity
\begin{equation}
 \boxed{
 H_{N,k}(\rho)=\rho^{k(N-1)}H_{N,k}(\rho^{-1}).}
 \label{eq:H-reciprocity}
\end{equation}
The Gaussian coefficient itself satisfies
\begin{equation}
 \qbinom{N}{k}{\rho}
 =\rho^{k(N-k)}\qbinom{N}{k}{\rho^{-1}}.
 \label{eq:qbinom-reciprocity}
\end{equation}

\subsection{Exact sign laws at negative nome}

For $0<a<1$, every factor $1-(-a)^j$ is positive, so
$\qbinom{N}{k}{-a}>0$.  Equation \eqref{eq:H-Gaussian} gives
\begin{equation}
 \boxed{
 \sgn H_{N,k}(-a)=(-1)^{\binom k2}
 \qquad(0<a<1).}
 \label{eq:H-sign-inside}
\end{equation}
For $a>1$, use reciprocity and the positive coefficient at $-1/a$:
\begin{equation}
 \boxed{
 \sgn H_{N,k}(-a)=
 (-1)^{\binom k2+k(N-k)}
 \qquad(a>1).}
 \label{eq:H-sign-outside}
\end{equation}
Thus the four signs in each requested orbit are controlled by two elementary
quadratic parities: the internal staircase $\binom k2$ and the rectangular
complement $k(N-k)$.

\subsection{Even/odd position splitting}

Negative nome separates even and odd positions.  Pairing the factors gives
\begin{align}
 (z;-a)_{2m}&=(z;a^2)_m(-az;a^2)_m,
 \label{eq:negative-Pochhammer-even}\\
 (z;-a)_{2m+1}&=(z;a^2)_{m+1}(-az;a^2)_m.
 \label{eq:negative-Pochhammer-odd}
\end{align}
This is the finite counterpart of the affine law for $Y_{-a}$ and the
decimation of $Y_q$ into $q^2$ components.  The negative Gaussian row is a
convolution of two positive $a^2$ rows, one carrying an extra sign and one
possibly an extra position.

\subsection{The finite-field/unitary square}

For a prime power $Q$, the positive Gaussian coefficient is the number of
$k$-planes in $\mathbb F_Q^N$:
\begin{equation}
 \qbinom{N}{k}{Q}=\#\Gr(k,\mathbb F_Q^N).
 \label{eq:Grassmann}
\end{equation}
Let $d=k(N-k)$.  Reciprocity gives
\begin{equation}
 \qbinom{N}{k}{1/Q}
 =Q^{-d}\#\Gr(k,\mathbb F_Q^N).
 \label{eq:Grassmann-reciprocal}
\end{equation}
Fu, Reiner, Stanton, and Thiem proved that negative Gaussian coefficients count
nondegenerate subspaces of a Hermitian space.  In their normalization,
\begin{equation}
 \qbinom{N}{k}{-Q}=(-Q)^{-d}\,U_Q(N,k),
 \label{eq:unitary-negative-q}
\end{equation}
where $U_Q(N,k)$ is the number of nondegenerate $k$-subspaces of a
nondegenerate Hermitian $N$-space over $\mathbb F_{Q^2}$.  Combining this with
reciprocity gives
\begin{equation}
 \qbinom{N}{k}{-1/Q}=Q^{-2d}U_Q(N,k).
 \label{eq:unitary-reciprocal}
\end{equation}

The requested parameters therefore form the square
\begin{center}
\begin{tabular}{c c c}
\toprule
nome & positive geometry & negative geometry \\
\midrule
$\pm2$, $\pm1/2$ & Grassmannians over $\mathbb F_2$ & Hermitian geometry over $\mathbb F_4$ \\
$\pm4$, $\pm1/4$ & Grassmannians over $\mathbb F_4$ & Hermitian geometry over $\mathbb F_{16}$ \\
\bottomrule
\end{tabular}
\end{center}
Multiplying by $\rho^{\binom k2}$ transports these counts to the
position-energy layers $H_{N,k}(\rho)$.  The file
\path{data/grassmannian_unitary_square.csv} records exact examples through
$N=7$ and checks both reciprocal normalizations with integer arithmetic.

\begin{example}
For $Q=2$, $N=4$, $k=2$, one has
\[
 \#\Gr(2,\mathbb F_2^4)=35,
 \qquad U_2(4,2)=240.
\]
Hence
\[
 \qbinom42{1/2}=\frac{35}{16},
 \qquad
 \qbinom42{-1/2}=\frac{15}{16}.
\]
The second rational number is $240/2^8$; the extra power relative to the
positive reciprocal reflects the unitary normalization $Q^{-2d}$.
\end{example}

\section[A geometric q-Prouhet moment-transfer theorem]{A geometric $q$-Prouhet moment-transfer theorem}
\label{sec:q-prouhet}

\subsection{Signed moments of geometric subset sums}

Define
\begin{equation}
 S_{N,m}(q)=
 \sum_{\epsilon\in\{0,1\}^N}
 (-1)^{|\epsilon|}A_{q,N}(\epsilon)^m.
 \label{eq:S-Nm}
\end{equation}
At $q=2$, these are the usual Thue--Morse signed power sums on the integer
block $0,1,\ldots,2^N-1$ after binary decoding.  At $q=1/2$ they are signed
moments of a finite dyadic subset grid.  At $q=1/4$ they live on a finite
base-$4$ Cantor skeleton.

The generating product is the Prouhet face \eqref{eq:Z-Prouhet-face}:
\begin{equation}
 \sum_{m\ge0}S_{N,m}(q)\frac{t^m}{m!}
 =\prod_{j=0}^{N-1}(1-\e^{q^jt}).
 \label{eq:S-EGF}
\end{equation}
Let
\begin{equation}
 T_{q,N}=\sum_{j=0}^{N-1}q^jU_j.
 \label{eq:T-qN}
\end{equation}
Its MGF is $\prod_{j<N}E(q^jt)$.

\begin{theorem}[Exact $q$-Prouhet moment transfer]
\label{thm:q-Prouhet}
For every nonzero complex $q$, every $N\ge1$, and every $\ell\ge0$,
\begin{align}
 S_{N,m}(q)&=0 &&(0\le m<N),
 \label{eq:q-Prouhet-vanish}\\
 S_{N,N+\ell}(q)
 &=\boxed{(-1)^N\frac{(N+\ell)!}{\ell!}
 q^{\binom N2}\E T_{q,N}^{\ell}.}
 \label{eq:q-Prouhet-transfer}
\end{align}
The theorem is \proved{} and \auditnew.
\end{theorem}

\begin{proof}
Factor each term on the right of \eqref{eq:S-EGF} as
$1-\e^{q^jt}=-q^jtE(q^jt)$.  Then
\begin{equation}
 \prod_{j=0}^{N-1}(1-\e^{q^jt})
 =(-1)^Nq^{\binom N2}t^N
  \prod_{j=0}^{N-1}E(q^jt).
 \label{eq:q-Prouhet-factor}
\end{equation}
The explicit factor $t^N$ proves the vanishings.  The coefficient of
$t^{N+\ell}$ on the right is
$(-1)^Nq^{\binom N2}\E T_{q,N}^{\ell}/\ell!$, whereas on the left it is
$S_{N,N+\ell}/(N+\ell)!$.
\end{proof}

The first surviving moment is
\begin{equation}
 S_{N,N}(q)=(-1)^NN!q^{\binom N2}.
 \label{eq:first-surviving}
\end{equation}
The theorem not only states an order of cancellation; it identifies every
post-cancellation coefficient with a positive-uniform moment when $q>0$.

\subsection[Recovering infinite q-Fabius moments]{Recovering infinite $q$-Fabius moments}

For $0<|q|<1$, $T_{q,N}$ converges in every $L^p$ to
$T_q=Y_q/(1-q)$.  Rearranging \eqref{eq:q-Prouhet-transfer} yields
\begin{equation}
 \boxed{
 m_\ell(q)=
 (1-q)^\ell\lim_{N\to\infty}
 (-1)^Nq^{-\binom N2}
 \frac{\ell!}{(N+\ell)!}S_{N,N+\ell}(q).}
 \label{eq:moment-from-Prouhet}
\end{equation}
Thus every Fabius or quartic moment is a normalized limit of the first
nonvanishing jets of finite Thue--Morse blocks.  This formula is the moment
version of the full MGF limit \eqref{eq:M-renormalized-Prouhet}.

The Python verification evaluates both sides of
\eqref{eq:q-Prouhet-transfer} exactly for all eight requested $q$ values,
$1\le N\le7$, and $0\le\ell\le4$.  All residuals in
\path{data/q_prouhet_verification.csv} are identically zero.

\subsection{Sign and reciprocal transport of the finite comb}

Introduce the signed atomic measure
\begin{equation}
 \Dcomb_{q,N}=\sum_{\epsilon}(-1)^{|\epsilon|}
 \delta_{A_{q,N}(\epsilon)}.
 \label{eq:Delta-qN}
\end{equation}
Its bilateral Laplace transform is \eqref{eq:S-EGF}.  Two elementary digit
bijections describe the full sign/reciprocal orbit.

For $q=-a$, $0<a<1$, complement the odd-position digits.  If
$m=\floor{N/2}$ and
\begin{equation}
 c_{a,N}=a\frac{1-a^{2m}}{1-a^2},
 \label{eq:c-aN}
\end{equation}
then
\begin{equation}
 \boxed{
 \Dcomb_{-a,N}=(-1)^m\,\tau_{-c_{a,N}}\Dcomb_{a,N},}
 \label{eq:comb-negative-translate}
\end{equation}
where $\tau_c$ translates a measure by $c$.  Digit reversal gives
\begin{equation}
 \boxed{
 \Dcomb_{q^{-1},N}
 =(D_{q^{-(N-1)}})_*\Dcomb_{q,N},}
 \label{eq:comb-reciprocal-dilate}
\end{equation}
where $D_c(x)=cx$.  Hence the finite Thue--Morse comb itself carries the same
Klein-four geometry as the infinite moment germ: sign is translation (up to a
global parity), inversion is dilation.

\section{Box-spline derivative combs and the quartic Cantor skeleton}
\label{sec:comb}

\subsection{Finite centered uniform sums}

For $0<q<1$ and $N\ge1$, let
\begin{equation}
 R_{q,N}=\sum_{j=0}^{N-1}q^jV_j,
 \qquad
 A_{q,N}^{\mathrm{rad}}=\frac12\sum_{j=0}^{N-1}q^j.
 \label{eq:RqN}
\end{equation}
The law of $R_{q,N}$ is a nonuniform univariate box spline.  Repeated
convolution of interval indicators gives the exact truncated-power formula.

\begin{proposition}[Finite box-spline formula]
\label{prop:box-spline}
The density $B_{q,N}$ of $R_{q,N}$ is
\begin{equation}
 \boxed{
 B_{q,N}(x)=
 \frac{1}{(N-1)!\,q^{\binom N2}}
 \sum_{\epsilon\in\{0,1\}^N}(-1)^{|\epsilon|}
 \pospart{x+A_{q,N}^{\mathrm{rad}}-A_{q,N}(\epsilon)}^{N-1}.}
 \label{eq:box-spline}
\end{equation}
Distributionally,
\begin{equation}
 \boxed{
 D^NB_{q,N}=
 q^{-\binom N2}
 \sum_{\epsilon}(-1)^{|\epsilon|}
 \delta_{-A_{q,N}^{\mathrm{rad}}+A_{q,N}(\epsilon)}.}
 \label{eq:box-spline-comb}
\end{equation}
The proposition is \proved{} and \derived.
\end{proposition}

\begin{proof}
The density of $q^jV_j$ is $q^{-j}$ times the indicator of
$[-q^j/2,q^j/2]$.  The standard inclusion--exclusion formula for a convolution
of $N$ interval densities with widths $q^j$ gives \eqref{eq:box-spline}; the
product of widths is $q^{\binom N2}$.  Since
$D^N\bigl((x-a)_+^{N-1}/(N-1)!\bigr)=\delta_a$, differentiating termwise proves
\eqref{eq:box-spline-comb}.
\end{proof}

The derivative comb is exactly the centered version of
$\Dcomb_{q,N}$.  Therefore the $q$-Prouhet cancellation theorem is also the
statement that its first $N-1$ distributional moments vanish.  Repeated
integration converts this signed atomic cancellation into a positive spline
density.

\subsection{Dyadic fullness versus quartic sparsity}

At $q=1/2$, the subset sums
$A_{1/2,N}(\epsilon)$ form a complete dyadic grid after a common scaling, and
the signs are the Thue--Morse sequence in binary order.  At $q=1/4$, the digits
occupy only the symbols $0$ and $1$ in base four.  In the infinite limit the
knot closure is the self-similar set
\begin{equation}
 K_4^{\mathrm{raw}}=\left\{\sum_{j\ge0}\epsilon_j4^{-j}:\epsilon_j\in\{0,1\}\right\},
 \qquad
 \dim_{\mathrm H}K_4^{\mathrm{raw}}=\frac{\log2}{\log4}=\frac12.
 \label{eq:quartic-Cantor}
\end{equation}
The finite derivative comb is therefore a signed Cantor skeleton, even though
the infinite convolution density itself is smooth and has full interval
support.

\begin{figure}[H]
\centering
\begin{subfigure}[t]{0.48\textwidth}
\centering
\includegraphics[width=\textwidth]{assets/dyadic_thue_morse_comb.png}
\caption{$q=1/2$: full dyadic grid.}
\end{subfigure}\hfill
\begin{subfigure}[t]{0.48\textwidth}
\centering
\includegraphics[width=\textwidth]{assets/quartic_cantor_comb.png}
\caption{$q=1/4$: base-$4$ digits $0,1$.}
\end{subfigure}
\caption{Signed derivative combs at level $N=9$.  The ordinate is the
Thue--Morse sign.  The quartic knots already display the dimension-$1/2$
Cantor geometry.}
\label{fig:combs}
\end{figure}

This contrast is conceptually important.  The derivative skeleton remembers
the discrete digit set, but integration through arbitrarily many uniform
scales erases the gaps.  Smoothness is not in conflict with a fractal hierarchy
of singular finite derivatives.

\subsection{Comb multisection}

Equation \eqref{eq:Z-multisection} at $(u,\rho)=(1,1)$ implies the convolution
law
\begin{equation}
 \boxed{
 \Dcomb_{q,mN}=*_{r=0}^{m-1}(D_{q^r})_*\Dcomb_{q^m,N}.}
 \label{eq:comb-multisection}
\end{equation}
For $m=2$,
\begin{equation}
 \Dcomb_{q,2N}=\Dcomb_{q^2,N}*(D_q)_*\Dcomb_{q^2,N}.
 \label{eq:comb-decimation}
\end{equation}
At $q=1/2$, the full dyadic Thue--Morse comb factors as a convolution of two
quartic Cantor combs.  This finite singular identity is the derivative-level
ancestor of the smooth distributional convolution
\eqref{eq:Fabius-quartic-convolution}.

\section{Exact inverse-geometric endpoint lattices}
\label{sec:endpoints}

\subsection{A simplex-volume identity}

Let $0<q<1$, let $G_q$ be the CDF of $Y_q$, and split the random series after
$n$ coordinates:
\begin{equation}
 Y_q=(1-q)\sum_{j=0}^{n-1}q^jU_j+q^nY_q',
 \label{eq:Y-tail-split}
\end{equation}
where $Y_q'$ is an independent copy of $Y_q$.  At arguments of the form
$q^nc$, the first $n$ coordinates are constrained to a simplex.  If the
threshold is below the smallest box width, none of the upper constraints
$U_j\le1$ is active.

\begin{theorem}[Stop-loss endpoint formula]
\label{thm:stop-loss-endpoint}
Let $n\ge1$ and
\begin{equation}
 0\le c\le\frac{1-q}{q}.
 \label{eq:c-range}
\end{equation}
Then
\begin{equation}
 \boxed{
 G_q(q^nc)=
 \frac{q^{n(n+1)/2}}{(1-q)^n n!}
 \E\pospart{c-Y_q}^{n}.}
 \label{eq:stop-loss-endpoint}
\end{equation}
The theorem is \proved{} and \auditnew.
\end{theorem}

\begin{proof}
Condition on $Y_q'=y$ in \eqref{eq:Y-tail-split}.  The event becomes
\[
 \sum_{j=0}^{n-1}(1-q)q^jU_j\le q^n(c-y).
\]
If $c-y<0$ its conditional probability is zero.  Otherwise the condition
\eqref{eq:c-range} implies
\[
 q^n(c-y)\le q^nc\le(1-q)q^{n-1},
\]
which is no larger than the smallest coordinate width.  Hence the intersection
with the unit cube is the full coordinate simplex, of volume
\[
 \frac{q^{n^2}(c-y)^n}
 {n!\prod_{j=0}^{n-1}(1-q)q^j}
 =\frac{q^{n(n+1)/2}}{(1-q)^nn!}(c-y)^n.
\]
Averaging over $y$ proves the formula.
\end{proof}

The right side is a normalized stop-loss moment.  The parameter $c$ therefore
interpolates an entire small-endpoint family, not merely a single lattice.
For $q>1/2$, the admissible interval ends before $c=1$; the general formula
still holds, but the particularly simple moment specialization below does not.

\subsection{The inverse-geometric moment identity}

For $0<q\le1/2$, $c=1$ is admissible.  Reflection of $Y_q$ gives
$1-Y_q\stackrel d=Y_q$, so \eqref{eq:stop-loss-endpoint} becomes the promised
generalization of the inverse-dyadic formula.

\begin{corollary}[Inverse-geometric endpoint]
\label{cor:inverse-geometric}
For $0<q\le1/2$ and $n\ge0$,
\begin{equation}
 \boxed{
 G_q(q^n)=
 \frac{q^{\binom{n+1}{2}}}{(1-q)^n}
 \frac{m_n(q)}{n!}.}
 \label{eq:inverse-geometric}
\end{equation}
Equivalently,
\begin{equation}
 \boxed{
 G_q(q^n)=
 \frac{q^{\binom{n+1}{2}}}{(q;q)_n}\mathcal P_n(q).}
 \label{eq:inverse-geometric-Pochhammer}
\end{equation}
This is \proved{} and \auditnew{} for general $q$.  The case $n=0$ is
immediate from $G_q(1)=1$; the preceding theorem supplies every case $n\ge1$.
\end{corollary}

At $q=1/2$, the prefactor simplifies to $2^{-\binom n2}$ and one recovers the
classical formula
\begin{equation}
 F(2^{-n})=2^{-\binom n2}\frac{m_n(1/2)}{n!}.
 \label{eq:classical-inverse-dyadic}
\end{equation}
At $q=1/4$, one obtains a new exact inverse-quartic lattice:
\begin{equation}
 G_{1/4}(4^{-n})
 =\frac{4^{-\binom{n+1}{2}}}{(3/4)^n}
 \frac{m_n(1/4)}{n!}.
 \label{eq:inverse-quartic}
\end{equation}
The first values are
\begin{equation}
 G_{1/4}(1/4)=\frac16,
 \quad
 G_{1/4}(1/16)=\frac1{240},
 \quad
 G_{1/4}(1/64)=\frac1{51840},
 \quad
 G_{1/4}(1/256)=\frac{121}{6768230400}.
 \label{eq:quartic-values}
\end{equation}
The table \path{data/inverse_geometric_endpoint_values.csv} contains exact
values through $n=12$ for both $q=1/2$ and $q=1/4$.

\begin{figure}[H]
\centering
\includegraphics[width=0.78\textwidth]{assets/quartic_endpoint_decay.png}
\caption{Exact inverse-quartic values plotted on a logarithmic ordinate.  The
quadratic growth of $-\log G_{1/4}(4^{-n})$ is the discrete shadow of the
endpoint asymptotic in \cref{sec:endpoint-asymptotic}.}
\label{fig:quartic-endpoint}
\end{figure}

\subsection{All derivative jets on the lattice}

Conditioning on the first coordinate gives the density equation
\begin{equation}
 G_q'(x)=\frac1{1-q}
 \left[
 G_q\!\left(\frac{x}{q}\right)
 -G_q\!\left(\frac{x-1+q}{q}\right)
 \right].
 \label{eq:CDF-density-equation}
\end{equation}
For $0\le x\le1-q$, the second term is zero.  Iterating at $x=q^n$ yields an
entire derivative jet.

\begin{theorem}[Endpoint derivative jet]
\label{thm:endpoint-jet}
Let $0<q\le1/2$, $n\ge0$, and $0\le r\le n$.  Then
\begin{equation}
 \boxed{
 G_q^{(r)}(q^n)=
 (1-q)^{-r}q^{-\binom r2}G_q(q^{n-r}).}
 \label{eq:endpoint-jet}
\end{equation}
Consequently
\begin{equation}
 G_q^{(r)}(q^n)=
 \frac{q^{\binom{n-r+1}{2}-\binom r2}}
 {(1-q)^n}
 \frac{m_{n-r}(q)}{(n-r)!}.
 \label{eq:endpoint-jet-moment}
\end{equation}
The theorem is \proved{} and \auditnew.
\end{theorem}

\begin{proof}
On the low interval, \eqref{eq:CDF-density-equation} reads
$G_q'(x)=(1-q)^{-1}G_q(x/q)$.  Differentiating $r-1$ more times and iterating
produces a factor
\[
 (1-q)^{-r}q^{-(0+1+\cdots+r-1)}.
\]
At $x=q^n$, all intermediate points remain in the low interval because
$q\le1/2$; at the boundary case $q=1/2$, the differentiated endpoint identities
follow by the flat matching of the two smooth branches.  This proves \eqref{eq:endpoint-jet}; substitute
\eqref{eq:inverse-geometric} for the second formula.
\end{proof}

For the density $g_q=G_q'$, the first case is
\begin{equation}
 g_q(q^n)=
 \frac{q^{n(n-1)/2}}{(1-q)^n}
 \frac{m_{n-1}(q)}{(n-1)!}.
 \label{eq:density-endpoint}
\end{equation}
At derivative order $r=n$, \eqref{eq:endpoint-jet} reduces to the elementary
value $(1-q)^{-n}q^{-\binom n2}$ because $G_q(1)=1$.

\subsection{Transport to negative parameters}

Let $G_{-a}$ denote the CDF of the negative law.  From
\cref{thm:negative-affine},
\begin{equation}
 G_{-a}(x)=G_a\!\left(\frac{x+\beta_a}{\lambda_a}\right).
 \label{eq:negative-CDF-affine}
\end{equation}
Therefore the positive endpoint lattice maps to
\begin{equation}
 x_{a,n}^{-}=\lambda_aa^n-\beta_a,
 \qquad
 \boxed{G_{-a}(x_{a,n}^{-})=G_a(a^n).}
 \label{eq:negative-endpoint-lattice}
\end{equation}
For $a=1/2$ these points are $3\cdot2^{-n}-1$ and converge to the left endpoint
$-1$.  For $a=1/4$ they are $(5/3)4^{-n}-1/3$ and converge to $-1/3$.
All derivative jets transfer by the chain rule with factor $\lambda_a^{-r}$.

\subsection{Successive endpoint ratios}

Since the support of $Y_q$ reaches $1$, its moment ratios satisfy
$m_{n+1}(q)/m_n(q)\to1$.  Dividing consecutive instances of
\eqref{eq:inverse-geometric} gives
\begin{equation}
 \frac{G_q(q^{n+1})}{G_q(q^n)}
 =\frac{q^{n+1}}{(1-q)(n+1)}
  \frac{m_{n+1}(q)}{m_n(q)}
 \sim\frac{q^{n+1}}{(1-q)(n+1)}.
 \label{eq:endpoint-ratio}
\end{equation}
This ratio is much smaller than any fixed geometric rate; it is the first
visible sign of the log-squared endpoint decay.

\section{Rigorous endpoint and inverse asymptotics}
\label{sec:endpoint-asymptotic}

\subsection{A uniform two-term logarithmic law}

The existing corpus contains a much sharper, all-orders, periodically corrected
small-$x$ expansion for the classical $q=1/2$ function.  Here we seek a result
that is weaker but uniform in every fixed $0<q\le1/2$, including $q=1/4$.
Put
\begin{equation}
 L=\log(1/q),
 \qquad
 \ell=\log(1/x).
 \label{eq:L-ell}
\end{equation}

Reflection and continuity imply
$\Prob(Y_q\ge1/2)=1/2$.  Hence
\begin{equation}
 2^{-(n+1)}\le m_n(q)\le1.
 \label{eq:moment-crude-bounds}
\end{equation}
These elementary bounds are sufficient to determine the first two growing
orders.

\begin{theorem}[Uniform two-term endpoint asymptotic]
\label{thm:endpoint-asymptotic}
Fix $0<q\le1/2$.  As $x\downarrow0$,
\begin{equation}
 \boxed{
 \log G_q(x)=
 -\frac{\ell^2}{2L}
 -\frac{\ell\log\ell}{L}
 +\bigO_q(\ell).}
 \label{eq:endpoint-asymptotic}
\end{equation}
The theorem is \proved{} and \auditnew{} for general $q$.
\end{theorem}

\begin{proof}
First put $x=q^n$, so $\ell=nL$.  Taking logarithms in
\eqref{eq:inverse-geometric} gives
\begin{equation}
 \log G_q(q^n)
 =-\frac{L}{2}n(n+1)-n\log(1-q)-\log(n!)+\log m_n(q).
 \label{eq:log-endpoint-n}
\end{equation}
By \eqref{eq:moment-crude-bounds}, $\log m_n(q)=O(n)$.  Stirling's formula
gives
$\log(n!)=n\log n-n+O(\log n)$.  Substituting $n=\ell/L$ into
\eqref{eq:log-endpoint-n} yields
\[
 \log G_q(q^n)
 =-\frac{\ell^2}{2L}
  -\frac{\ell\log\ell}{L}+O_q(\ell).
\]

For arbitrary $x$, choose $n$ such that $q^{n+1}<x\le q^n$.  Monotonicity
places $G_q(x)$ between the two adjacent lattice values.  Replacing $\ell$ by
$nL$ or $(n+1)L$ changes the two displayed leading terms by $O_q(\ell)$, so the
same expansion holds throughout the interval.
\end{proof}

The leading term alone gives the familiar log-normal-like flatness.  The
second term is larger than the unspecified linear correction and is therefore
a genuinely stable feature of the whole geometric family.

For the two principal cases,
\begin{align}
 \log F(x)
 &=-\frac{\log^2(1/x)}{2\log2}
   -\frac{\log(1/x)\log\log(1/x)}{\log2}
   +O(\log(1/x)),
 \label{eq:dyadic-two-term}\\
 \log G_{1/4}(x)
 &=-\frac{\log^2(1/x)}{4\log2}
   -\frac{\log(1/x)\log\log(1/x)}{2\log2}
   +O(\log(1/x)).
 \label{eq:quartic-two-term}
\end{align}
The quartic law is less flat at a common numerical $x$ because its scale
ladder descends more rapidly.

\subsection{Asymptotic inversion}

Let $y=G_q(x)$ and $S=\log(1/y)$.  Equation
\eqref{eq:endpoint-asymptotic} reads
\begin{equation}
 S=\frac{\ell^2}{2L}+\frac{\ell\log\ell}{L}+O_q(\ell).
 \label{eq:S-ell}
\end{equation}
The first approximation is $\ell\sim\sqrt{2LS}$.  One substitution then
captures the logarithmic correction.

\begin{corollary}[Inverse endpoint asymptotic]
\label{cor:inverse-asymptotic}
For fixed $0<q\le1/2$,
\begin{equation}
 \boxed{
 \log\frac1{G_q^{-1}(y)}
 =\sqrt{2L\log(1/y)}
 -\frac12\log\log(1/y)+O_q(1),
 \qquad y\downarrow0.}
 \label{eq:inverse-asymptotic}
\end{equation}
This is \proved.
\end{corollary}

\begin{proof}
Let $A=\sqrt{2LS}$.  Set $\ell=A+\delta$.  Inserting this into
\eqref{eq:S-ell}, the terms of size $A\log A$ are
\[
 \frac{A}{L}(\delta+\log A).
\]
Thus $\delta=-\log A+O_q(1)$.  Since
$\log A=\frac12\log S+O_q(1)$, the stated formula follows.  Standard
monotonicity of the left side of \eqref{eq:S-ell} controls the $O(1)$
remainder.
\end{proof}

For negative $q=-a$, apply the affine inverse
\eqref{eq:negative-CDF-affine}.  Near the left endpoint
$\ell_a=-a/(1-a)$,
\begin{equation}
 G_{-a}^{-1}(y)-\ell_a
 =\lambda_aG_a^{-1}(y),
 \label{eq:negative-inverse-transfer}
\end{equation}
so the same square-root/log-log law holds after replacing distance to zero by
distance to the left endpoint.

\subsection[A Lambert-W-minus-one saddle coordinate]{A Lambert-$W_{-1}$ saddle coordinate}

The all-orders classical analysis in the corpus uses the lower real branch of
Lambert $W$.  The same coordinate is predicted for every $q$, but proving the
periodic correction and the complete saddle expansion uniformly in $q$ is a
separate project.

Let $\mathcal L_q(s)=M_q(-s)=\E\e^{-sY_q}$ for $0<q<1$.  From the product
or from Mellin analysis of the scale sum, one expects the saddle equation
\begin{equation}
 x\sim\frac{\log s}{s\log(1/q)}.
 \label{eq:saddle-equation}
\end{equation}
Solving its dominant part gives
\begin{equation}
 s\sim-\frac1{x\log(1/q)}
 W_{-1}(-C_qx),
 \label{eq:Lambert-saddle}
\end{equation}
where $C_q>0$ depends on the normalization chosen for the lower-order terms.
The log-periodic fluctuation should have period one in
$\log_{1/q}s$.

\begin{conjecture}[General geometric all-orders endpoint expansion]
\label{conj:all-orders-q}
For each $0<q<1$, the density $g_q(x)$ admits as $x\downarrow0$ a complete
saddle expansion in inverse powers of
$-W_{-1}(-C_qx)$, with analytic coefficient functions periodic in
$\log_{1/q}(-W_{-1}(-C_qx)/x)$.  The leading two unbounded terms agree with
\cref{thm:endpoint-asymptotic}.  The conjecture is \openstatus{} for general
$q$; the classical $q=1/2$ case has substantially stronger support in the
existing corpus.
\end{conjecture}

The quartic case is especially attractive.  Scale decimation
\eqref{eq:Fabius-quartic-convolution} couples its periodic correction to the
classical one, but at two interlaced scales.  This may allow the quartic
periodic function to be reconstructed from the dyadic function through a
nonlinear convolution saddle equation rather than recomputed from scratch.

\section{Numerical experiments and reproducibility}
\label{sec:numerics}

\subsection{Scope of the computation}

The accompanying program \path{numerical_experiments.py} is intentionally
self-contained.  Exact calculations use Python's \texttt{fractions.Fraction}
and integer arithmetic.  The only floating-point components are independent
high-precision product checks and deterministic Fourier inversions for the
figures.  No random sampling is used.

The script verifies:
\begin{enumerate}
\item the moment recurrence through order $10$ at
      $q=\pm1/2,\pm2,\pm1/4,\pm4$;
\item negative-$q$ affine conjugacy at every checked moment;
\item reciprocal binomial-convolution inversion through order $10$;
\item $q\mapsto q^2$ decimation for $q=\pm1/2,\pm2$;
\item the position-energy identity
      $H_{N,k}(q)=q^{\binom k2}\qbinom Nkq$ for $N\le8$;
\item the $q$-Prouhet transfer for $N\le7$, $\ell\le4$, at all eight values;
\item the exact moment-polynomial recurrence and the conjectural odd-$q$-integer
      divisor by rational polynomial long division through $n=16$;
\item inverse-geometric endpoints through $n=12$ for $q=1/2,1/4$;
\item the Grassmannian/unitary normalizations for $Q=2,4$ and $N\le7$;
\item the reciprocal Hankel signature through order $9$ for
      $q=\pm2,\pm4$;
\item the Mahler, sign, reciprocal, and decimation product identities at
      $90$ decimal digits.
\end{enumerate}

The largest high-precision residual is approximately $1.3\times10^{-40}$; the
remaining analytic residuals are typically around $10^{-90}$.  The weaker
sign-conjugacy residual reflects the deliberately finite $300$-factor
truncation after an argument dilation as large as five, not a loss of the exact
identity.

\subsection{Generated files}

\begin{center}
\begin{tabularx}{0.96\textwidth}{>{\ttfamily}l X}
\toprule
\textnormal{file} & \textnormal{content} \\
\midrule
\path{analytic_product_errors.csv} & high-precision residuals for four analytic identities \\
\path{q_orbit_moments.csv} & exact moments and variances for all eight requested parameters \\
\path{q_prouhet_verification.csv} & exact two-sided values in \cref{thm:q-Prouhet} \\
\path{moment_polynomial_divisibility.csv} & exact $\mathcal P_n$ coefficients, divisors, and quotients through $n=16$ \\
\path{inverse_geometric_endpoint_values.csv} & exact dyadic and quartic endpoint values \\
\path{grassmannian_unitary_square.csv} & positive/negative finite-geometry normalizations \\
\path{reciprocal_hankel_signatures.csv} & exact determinants and predicted signs \\
\bottomrule
\end{tabularx}
\end{center}

The program can be rerun from the report directory with
\begin{lstlisting}[style=python]
python numerical_experiments.py
\end{lstlisting}
It creates the \path{data/} and \path{assets/} directories if needed, rewrites
the CSV files, verifies all assertions, and regenerates the five figures.

\subsection{Selected exact moment rows}

For reference, the first six raw moments are shown below.  They make the
reciprocal sign change tangible: positive contractions have positive variance,
whereas their reciprocal germs have the opposite centered cumulants.

\begin{table}[H]
\centering
\small
\caption{Exact moments $m_n(q)$ for the requested parameter values.}
\label{tab:moments}
\begin{tabular}{c c c c c c c}
\toprule
$q$ & $m_1$ & $m_2$ & $m_3$ & $m_4$ & $m_5$ & $m_6$ \\
\midrule
$1/2$ & $1/2$ & $5/18$ & $1/6$ & $143/1350$ & $19/270$ & $1153/23814$ \\
$-1/2$ & $1/2$ & $1/2$ & $1/2$ & $29/50$ & $7/10$ & $439/490$ \\
$2$ & $1/2$ & $2/9$ & $1/12$ & $16/675$ & $1/270$ & $-16/59535$ \\
$-2$ & $1/2$ & $0$ & $-1/4$ & $-2/25$ & $3/10$ & $10/49$ \\
\addlinespace
$1/4$ & $1/2$ & $3/10$ & $1/5$ & $121/850$ & $9/85$ & $44017/541450$ \\
$-1/4$ & $1/2$ & $7/18$ & $1/3$ & $283/918$ & $137/459$ & $1573771/5262894$ \\
$4$ & $1/2$ & $1/5$ & $1/20$ & $-1/425$ & $-1/170$ & $121/54145$ \\
$-4$ & $1/2$ & $1/9$ & $-1/12$ & $-31/459$ & $49/918$ & $213523/2631447$ \\
\bottomrule
\end{tabular}
\end{table}

\subsection{Independent checks versus proof}

The exact computations were also used as a theorem-discovery filter.  Several
plausible but false reciprocity formulas for the pole-cleared polynomials
$\mathcal P_n$ were rejected before inclusion.  Only the recurrence,
boundaries, and endpoint identity proved in \cref{sec:moment-polynomial} are
stated as theorems.  The cyclotomic divisor remains a conjecture despite its
symbolic evidence.  This separation is important in a subject where
$q\leftrightarrow1/q$ often introduces triangular powers and can make an
incorrect formula look valid in low degree.

\section{Conjectures and frontier directions}
\label{sec:conjectures}

\subsection{Natural boundaries in the parameter plane}

The centered logarithm has the formal expansion
\begin{equation}
 \log C_q(t)=
 \sum_{m=1}^{\infty}
 \frac{B_{2m}}{2m}
 \frac{(1-q)^{2m}}{1-q^{2m}}
 \frac{t^{2m}}{(2m)!}.
 \label{eq:logC-q}
\end{equation}
For fixed nonzero $t$, the coefficient functions have poles at roots of unity
of arbitrarily high order.  Some singularities may cancel after exponentiation
or analytic continuation, and the point $q=1$ is a degenerate concentration
limit, so the dense-pole heuristic is not by itself a proof.

\begin{conjecture}[Generic natural boundary]
\label{conj:natural-boundary}
For generic fixed $t\ne0$, the unit circle is a natural boundary in the
$q$-plane for $C_q(t)$ after removal of any isolated removable degeneracies.
Equivalently, no single-valued meromorphic continuation crosses a nontrivial
arc of $|q|=1$.  This is \openstatus.
\end{conjecture}

A viable proof strategy is to isolate primitive root-of-unity singularities in
one cumulant order whose principal part cannot be cancelled by lower orders,
then control the exponential remainder.  The two-nome partition function may
supply finite approximants with explicit root-of-unity zero orders.

\subsection{Arithmetic of the moment polynomials}

The polynomials $\mathcal P_n(q)$ connect moments, Bernoulli limits,
$q$-Pochhammer denominators, and exact endpoint values.  Besides
\cref{conj:Pn-divisibility}, several questions are immediate.

\begin{problem}[Coefficient signs and unimodality]
After dividing $\mathcal P_n$ by the conjectural odd-$q$-integer factor, do the
remaining coefficients have a stable sign pattern, log-concavity, or
unimodality?  The first factorizations suggest that parity and Bernoulli signs
must be separated before any positivity statement is formulated.
\end{problem}

\begin{problem}[Root geometry]
Describe the zeros of $\mathcal P_n$.  Which are forced roots of unity, which
lie on reciprocal curves, and do the noncyclotomic zeros approach a limiting
set?  Is there a transform relating them to the vertical zeros of the
reciprocal orthogonal polynomials in \cref{sec:laplace-dual}?
\end{problem}

\begin{problem}[Denominator theorem]
Find the least integer $D_n$ such that $D_n\mathcal P_n(q)\in\Z[q]$.  Compare
its prime valuations with the von Staudt--Clausen denominators of $B_{2m}$ and
the Mersenne/Fermat-type factors in $(q;q)_n$ at rational $q$.
\end{problem}

\subsection{The reciprocal orthogonal-polynomial problem}

The positive measure $L_r$ is an infinite convolution of Laplace laws at the
double geometric/harmonic scales
\begin{equation}
 \frac{|1-r||r|^j}{2\pi k}.
 \label{eq:double-scales}
\end{equation}
It is symmetric, analytic, and infinitely divisible.  Its orthogonal
polynomials are therefore a well-posed new family.  Important questions
include:
\begin{enumerate}
\item determine the Jacobi coefficients $\beta_n^{(r)}$ exactly or
      asymptotically;
\item find a $q$-difference or integro-difference equation induced by adding
      one geometric layer of Laplace variables;
\item determine the limiting zero distribution after the natural scaling;
\item compare $r=1/2$ and $r=1/4$ through scale multisection;
\item identify whether the laws belong to a tractable subclass of generalized
      gamma convolutions and exploit the corresponding Thorin measure.
\end{enumerate}
The last point should be handled carefully: individual Laplace variables are
differences of exponentials, but the signed difference structure is not itself
a positive gamma convolution on $\R_+$.

\subsection{A two-nome asymptotic phase diagram}

The master function $\Zpart_N(u,t;\rho,\sigma)$ invites asymptotics in which
$N\to\infty$ while $k/N$, $m/N$, or the nomes vary.  Possible regimes include:
\begin{itemize}
\item $|\rho|<1$, $|\sigma|<1$: convergent position and projection energies;
\item $\rho>1$, $\sigma>1$: large-deviation concentration on terminal digits;
\item $\rho$ or $\sigma$ on the unit circle: resonance and partial Prouhet
      cancellation;
\item reciprocal pairs $(\rho,\sigma)\mapsto(1/\rho,1/\sigma)$: digit reversal
      and dilation;
\item negative pairs: parity splitting and affine translation of the projected
      comb.
\end{itemize}
A saddle-point analysis of the mixed coefficients
$\mathcal M_{N,k,m}$ could connect Gaussian-binomial limit shapes to
Thue--Morse rarefied sums.  This appears to be a genuinely two-dimensional
problem; specializing $\rho=\sigma$ too early loses the phase diagram.

\subsection{Quartic Cantor renormalization}

At $q=1/4$, finite derivative knots lie on a Cantor skeleton but the limiting
density is $C^\infty$ on a full interval.  The exact dyadic--quartic
factorization provides a renormalization relation between a full grid and two
sparse grids.  Promising problems are:
\begin{enumerate}
\item quantify how the Hausdorff geometry of the knots controls finite-level
      derivative norms;
\item determine the best uniform and Sobolev convergence rates of quartic box
      splines to the limiting density;
\item derive a wavelet or multiresolution basis adapted to the digit set
      $\{0,1\}$ in base four;
\item compare the zero multiplicity divisor
      $1+\floor{\nu_2(n)/2}$ with the dyadic divisor $1+\nu_2(n)$ under the
      convolution \eqref{eq:up-quartic-convolution};
\item use the exact endpoint lattice to calibrate numerical inverse-CDF
      algorithms near zero.
\end{enumerate}

\section{Lean formalization roadmap}
\label{sec:lean-roadmap}

The existing formal development already contains the generic weighted-uniform
series, geometric weights, probability law, affine self-similarity, reflection,
support, absolute continuity, and the classical Fabius bridge.  The results of
this report can therefore be added incrementally rather than by introducing a
parallel foundation.

\subsection{Low-risk algebraic tranche}

The following statements require only finite sums, formal power series, or
existing probability transport lemmas.
\begin{enumerate}
\item \textbf{Negative affine conjugacy.}  Center each uniform coordinate,
      use invariance under coordinatewise reflection, and push forward by the
      affine map.
\item \textbf{Geometric multisection.}  Reindex the series by
      $n=mk+r$ and apply independence of coordinate subsequences.
\item \textbf{Moment recurrence and reciprocity.}  Define a formal or analytic
      germ interface and prove coefficient uniqueness away from roots of unity.
\item \textbf{Position-energy layer.}  Reuse the finite $q$-binomial core and
      prove the subset-energy expansion directly.
\item \textbf{Two-nome partition function.}  A finite-product definition and
      extensionality proof should make multisection nearly automatic.
\item \textbf{$q$-Prouhet transfer.}  Factor each exponential difference and
      compare coefficients in a finite exponential generating function.
\end{enumerate}
These modules would have small analytic dependencies and should compile
quickly under the repository's single-worker build policy.

\subsection{Endpoint-analysis tranche}

The endpoint theorem needs the finite-dimensional uniform-simplex volume.  A
formal proof can proceed by an affine change of variables from the standard
simplex
\[
 \{u_j\ge0:\sum u_j\le h\}
\]
and a lemma showing that the upper cube constraints are inactive under
\eqref{eq:c-range}.  Once \eqref{eq:stop-loss-endpoint} is formalized,
reflection gives the moment formula, and the derivative jet follows from the
existing CDF density equation.  The two-term asymptotic needs only Stirling,
moment bounds, and monotonic bracketing; it is far lighter than the existing
all-orders saddle machinery.

Suggested files are:
\begin{center}
\begin{tabular}{l l}
\toprule
module & principal content \\
\midrule
\texttt{GeometricUniformSign.lean} & negative affine law and support endpoints \\
\texttt{GeometricUniformMultisection.lean} & $q\mapsto q^m$ independent decomposition \\
\texttt{GeometricProuhet.lean} & finite transfer and normalized moment limit \\
\texttt{TwoNomePartition.lean} & master product and mixed coefficients \\
\texttt{GeometricEndpoint.lean} & stop-loss formula and exact lattice \\
\texttt{GeometricEndpointJets.lean} & derivative hierarchy \\
\texttt{GeometricEndpointAsymptotic.lean} & two-term CDF and inverse laws \\
\bottomrule
\end{tabular}
\end{center}

\subsection{Reciprocal analytic tranche}

The reciprocal Laplace theorem requires more infrastructure:
\begin{enumerate}
\item define the reciprocal germ as a locally normally convergent inverse
      product;
\item prove the Euler product rearrangement;
\item construct the double-indexed Laplace series in $L^2$;
\item identify its characteristic function by finite partial products and
      dominated convergence;
\item establish strict positivity of real Hankel Gram determinants;
\item transport determinants under affine complex maps;
\item define the induced quasi-definite moment functional and orthogonal
      polynomial recurrence.
\end{enumerate}
A pragmatic split is to formalize reciprocity and cumulants first, then the
Laplace representation, and only then the Hankel/orthogonality consequences.

\subsection{Status separation in documentation}

The report's theorem/conjecture boundary should be preserved if material is
merged into the repository.  In particular, the following must remain in the
frontier quarantine until proved:
\begin{itemize}
\item odd-$q$-integer divisibility of $\mathcal P_n$;
\item generic natural boundary at $|q|=1$;
\item all-orders Lambert-$W_{-1}$ expansion for general $q$;
\item asymptotic statements about reciprocal Jacobi coefficients;
\item coefficient positivity or root-location patterns for $\mathcal P_n$.
\end{itemize}

\section{Conclusion}
\label{sec:conclusion}

The main structural lesson is that ``other $q$ values'' do not merely replace
$1/2$ by another number in an existing formula.  Three transformations act on
three distinct layers:
\begin{itemize}
\item $q\mapsto-q$ changes the outer affine normalization while leaving the
      centered Rvachev law invariant in distribution;
\item $q\mapsto1/q$ turns geometric smoothing into deconvolution, zeros into
      poles, and real probability moments into vertical-line complex moments;
\item $q\mapsto q^m$ separates scales into congruence classes and connects the
      dyadic and quartic systems by exact convolution.
\end{itemize}
The finite Gaussian nome and the geometric projection need not coincide.  The
two-nome partition function shows that $q$-Pochhammer coefficients and
Thue--Morse Prouhet moments are complementary statistics of the same binary
cube.  The normalized residual of its maximal Prouhet zero is exactly the
$q$-Fabius MGF.

For the requested values, the picture is now compact.  The orbit
$\{\pm1/2,\pm2\}$ is controlled by the classical Fabius law and finite geometry
over $\mathbb F_2/\mathbb F_4$; the orbit $\{\pm1/4,\pm4\}$ is controlled by a
quartic smooth law whose finite derivatives sit on a dimension-$1/2$ Cantor
skeleton and whose Gaussian coefficients correspond to
$\mathbb F_4/\mathbb F_{16}$ geometry.  Each positive contraction has an exact
inverse-geometric endpoint lattice, each negative contraction inherits it by
affine transport, and each reciprocal germ has a positive Laplace spectral
dual.

The strongest immediately actionable next steps are the Lean formalization of
negative affine conjugacy, multisection, the finite $q$-Prouhet transfer, and
the endpoint simplex theorem.  The most mathematically open directions are the
root-of-unity arithmetic of $\mathcal P_n$, the reciprocal vertical-line
orthogonal polynomials, and the all-orders quartic/general-$q$ Lambert-$W$
endpoint expansion.

\appendix

\section{Detailed identities for the requested parameter values}
\label{app:requested-values}

\subsection{Orbit transform table}

Let $M=M_{1/2}$ and $N=M_{1/4}$.  Then
\begin{center}
\begin{tabularx}{0.96\textwidth}{c X}
\toprule
parameter & exact germ in terms of the positive contraction \\
\midrule
$1/2$ & $M_{1/2}(t)=M(t)$ \\
$-1/2$ & $M_{-1/2}(t)=\e^{-t}M(3t)$ \\
$2$ & $M_2(t)=\e^t/M(t)$ \\
$-2$ & $M_{-2}(t)=\e^{2t}/M(3t)$ \\
\addlinespace
$1/4$ & $M_{1/4}(t)=N(t)$ \\
$-1/4$ & $M_{-1/4}(t)=\e^{-t/3}N(5t/3)$ \\
$4$ & $M_4(t)=\e^t/N(t)$ \\
$-4$ & $M_{-4}(t)=\e^{4t/3}/N(5t/3)$ \\
\bottomrule
\end{tabularx}
\end{center}

The centered transforms remove the exponential factors:
\begin{align*}
 C_{-1/2}(t)&=C_{1/2}(3t),
 &C_2(t)&=C_{1/2}(t)^{-1},
 &C_{-2}(t)&=C_{1/2}(3t)^{-1},\\
 C_{-1/4}(t)&=C_{1/4}(5t/3),
 &C_4(t)&=C_{1/4}(t)^{-1},
 &C_{-4}(t)&=C_{1/4}(5t/3)^{-1}.
\end{align*}

\subsection{Cumulant scaling table}

For $m\ge1$, let $K_{2m}^{(2)}=\kappa_{2m}(1/2)$ and
$K_{2m}^{(4)}=\kappa_{2m}(1/4)$.  Then
\begin{center}
\begin{tabular}{c c c c c}
\toprule
orbit & $q$ & $-q$ & $1/q$ & $-1/q$ \\
\midrule
half & $K_{2m}^{(2)}$ & $3^{2m}K_{2m}^{(2)}$ & $-K_{2m}^{(2)}$ & $-3^{2m}K_{2m}^{(2)}$ \\
quarter & $K_{2m}^{(4)}$ & $(5/3)^{2m}K_{2m}^{(4)}$ & $-K_{2m}^{(4)}$ & $-(5/3)^{2m}K_{2m}^{(4)}$ \\
\bottomrule
\end{tabular}
\end{center}
This table follows at once from affine scaling and reciprocal inversion.  It
is often more efficient than re-evaluating the rational formula
\eqref{eq:cumulants} eight times.

\section{Proof details for finite-comb transport}
\label{app:comb-proofs}

\subsection[Negative-q translation]{Negative-$q$ translation}

Let $q=-a$ and write
\begin{align*}
 A_{-a,N}(\epsilon)
 &=\sum_{j\ \mathrm{even}}\epsilon_ja^j
   -\sum_{j\ \mathrm{odd}}\epsilon_ja^j.
\end{align*}
Define $\epsilon'$ by preserving even digits and complementing odd digits.
Then
\begin{align*}
 A_{a,N}(\epsilon')
 &=\sum_{j\ \mathrm{even}}\epsilon_ja^j
   +\sum_{j\ \mathrm{odd}}(1-\epsilon_j)a^j\\
 &=A_{-a,N}(\epsilon)+c_{a,N}.
\end{align*}
If there are $m=\floor{N/2}$ odd positions, then
$(-1)^{|\epsilon|}=(-1)^m(-1)^{|\epsilon'|}$.  Summing the atoms proves
\eqref{eq:comb-negative-translate}.

\subsection{Reciprocal dilation}

Reverse the digits, $\epsilon'_j=\epsilon_{N-1-j}$.  Then
\begin{align*}
 A_{q^{-1},N}(\epsilon)
 &=\sum_{j=0}^{N-1}\epsilon_jq^{-j}\\
 &=q^{-(N-1)}\sum_{j=0}^{N-1}\epsilon_{N-1-j}q^j
 =q^{-(N-1)}A_{q,N}(\epsilon').
\end{align*}
Digit reversal preserves Hamming parity, proving
\eqref{eq:comb-reciprocal-dilate}.

\section{Moment-polynomial data and the divisor conjecture}
\label{app:Pn-data}

The first additional factorizations are
\begin{align}
 \mathcal P_5(q)
 &=-\frac{[3]_q[5]_q
 (q^3-2q^2-2)}{1440},
 \label{eq:P5-factor}\\
 \mathcal P_6(q)
 &=\frac{[3]_q[5]_q}{60480}
 \left(
 2q^9-6q^8-3q^7+4q^6+6q^5+15q^4
 +3q^3+24q^2+6q+12
 \right),
 \label{eq:P6-factor}\\
 \mathcal P_7(q)
 &=\frac{[3]_q[5]_q[7]_q}{120960}
 \left(
 2q^8-6q^7+4q^6-3q^5+6q^4-4q^3+7q^2+3
 \right).
 \label{eq:P7-factor}
\end{align}
Here $[d]_q=(q^d-1)/(q-1)$.  These formulas were obtained from the proved
recurrence \eqref{eq:Pn-recurrence}; they are exact.  The extrapolation to all
odd $d\le n$ is the conjectural step.

A root-of-unity proof would need to establish the correct multiplicity of each
cyclotomic polynomial.  Since
\[
 \prod_{\substack{3\le d\le n\\d\ \mathrm{odd}}}[d]_q
 =\prod_{\substack{3\le e\le n\\e\ \mathrm{odd}}}
 \Phi_e(q)^{\left\lfloor(\lfloor n/e\rfloor+1)/2\right\rfloor},
\]
the exponent of $\Phi_e$ counts the odd multiples of $e$ not exceeding $n$.
Thus simple vanishing of $\mathcal P_n$ at primitive roots is not enough.  One
must control repeated jets.  This is where the finite cancellation-order proposition
\ref{prop:Z-order} may become useful.

\section{Reproducibility notes}
\label{app:reproducibility}

The script requires Python 3.10 or newer and the packages
\texttt{numpy}, \texttt{matplotlib}, and \texttt{mpmath}.  Exact arithmetic
uses only the standard library.  The figures are deterministic:
\begin{itemize}
\item densities are Fourier inversions of truncated exact sinc products on a
      periodic grid of size $2^{16}$;
\item negative densities are produced by the exact affine change of variables;
\item comb figures enumerate all $2^9$ binary vectors;
\item endpoint decay uses exact rational moments before conversion to the
      plotted logarithm.
\end{itemize}
The source is extensively commented and each CSV row contains enough input
parameters to reproduce the underlying assertion independently.

\section{Bibliographic and corpus references}
\label{app:bibliography}

\begin{thebibliography}{99}

\bibitem{AriasDefinition}
J. Arias de Reyna,
\emph{An infinitely differentiable function with compact support: definition
and properties},
2017, arXiv:1702.05442.

\bibitem{AriasArithmetic}
J. Arias de Reyna,
\emph{Arithmetic of the Fabius function},
2017, arXiv:1702.06487, version 3.

\bibitem{FabiusOriginal}
J. Fabius,
\emph{A probabilistic example of a nowhere analytic $C^\infty$-function},
Zeitschrift f\"ur Wahrscheinlichkeitstheorie und Verwandte Gebiete
\textbf{5} (1966), 173--174.

\bibitem{Rvachev}
V. L. Rvachev,
work on compactly supported infinitely differentiable atomic functions and the
up-function; see the historical sources and bibliography collected in the
\repo{} primary exposition.

\bibitem{NegativeQ}
A. M. Fu, V. Reiner, D. Stanton, and P. Thiem,
\emph{The negative $q$-binomial},
Electronic Journal of Combinatorics \textbf{19} (2012), Paper P36;
arXiv:1108.4702.

\bibitem{DLMFq}
NIST Digital Library of Mathematical Functions,
Chapter 17, \emph{$q$-Hypergeometric and Related Functions}, especially the
$q$-Pochhammer and Gaussian-binomial conventions.

\bibitem{DLMFW}
NIST Digital Library of Mathematical Functions,
Section 4.13, \emph{Lambert $W$-Function}.

\bibitem{Andrews}
G. E. Andrews,
\emph{The Theory of Partitions},
Cambridge University Press, 1998 reprint.

\bibitem{GasperRahman}
G. Gasper and M. Rahman,
\emph{Basic Hypergeometric Series}, 2nd ed.,
Cambridge University Press, 2004.

\bibitem{StanleyEC1}
R. P. Stanley,
\emph{Enumerative Combinatorics, Volume 1}, 2nd ed.,
Cambridge University Press, 2012.

\bibitem{ProveItPrimary}
V. Reshetnikov,
\repo{} repository,
\path{Analysis/FabiusFunction/docs/Fabius_Function_and_Rvachev_Up/},
primary proof-backed exposition and Lean development.

\bibitem{ProveItFrontier}
V. Reshetnikov,
\repo{} repository,
\path{Analysis/FabiusFunction/docs/non-formalized-research-frontiers/},
canonical frontier volume and thematic syntheses, pinned at commit
\texttt{89b1a5f915cad065a42ee48f758eb2cd2d9beb94}.

\bibitem{ProveItQ}
V. Reshetnikov,
\repo{} archived expositions
\path{archive/q-calculus/Fabius_q_Special_Functions/} and
\path{archive/research-frontiers/Fabius_Dyadic_q_Connections/},
for the half-base $q$-Pochhammer/Gaussian framework and exact dyadic formulae.

\end{thebibliography}

\end{document}
